\input{estilo.tex}
\begin{document}
\begin{titlepage}\thispagestyle{empty}
\begin{tikzpicture}[remember picture,overlay]
\fill[Ink] (current page.south west) rectangle (current page.north east);
\fill[Teal] ([yshift=-14mm]current page.north west) rectangle ([yshift=-17mm]current page.north east);
\node[anchor=north west,align=left,text width=168mm,text=white] at ([xshift=21mm,yshift=-30mm]current page.north west) {
{\sffamily\fontsize{10}{13}\selectfont\bfseries APRENDER HACIENDO · MÓDULO OPCIONAL}\par\vspace{10mm}
{\sffamily\fontsize{34}{39}\selectfont\bfseries De radiografías\\a predicciones}\par\vspace{6mm}
{\sffamily\fontsize{17}{23}\selectfont Primer modelo de\\aprendizaje automático}\par\vspace{7mm}
{\sffamily\fontsize{11}{16}\selectfont Una guía visual de datos, entrenamiento,\\evaluación y errores, paso a paso.}};
\node[anchor=center,inner sep=4mm,fill=white,rounded corners=3pt] at ([yshift=-26mm]current page.center) {\includegraphics[width=160mm]{figuras/celda-01.png}};
\node[anchor=south west,text width=168mm,align=left,text=white] at ([xshift=21mm,yshift=23mm]current page.south west) {
{\sffamily\fontsize{13}{17}\selectfont\bfseries Juan Calderón}\par\vspace{3mm}
{\sffamily\fontsize{10}{15}\selectfont Machine Learning Honduras · 2026\\PneumoniaMNIST · Python · scikit-learn · CPU}\par\vspace{5mm}
{\sffamily\fontsize{9}{13}\selectfont Imágenes públicas para aprender. Este ejercicio no es una herramienta diagnóstica.}};
\end{tikzpicture}\null
\end{titlepage}
\setcounter{page}{2}



\Pagina{APRENDER A PARTIR DE EJEMPLOS}{De píxeles a una predicción}{Taller impartido el jueves 26 de noviembre de 2026 a las 13:30, hora de Honduras (UTC−6).}

Mira las radiografías de la portada y sus etiquetas. El programa recibirá los píxeles de otra imagen y propondrá una de esas dos etiquetas. Para aprender, dispone de ejemplos con su respuesta de referencia.

\href{https://colab.research.google.com/github/FisikMedik/taller-imagenes-medicas-python/blob/main/Modulo_ML_radiografias.ipynb}{Abrir el cuaderno del módulo en Google Colab}.

\Sub{Qué vamos a construir}

\textbf{Una imagen $\rightarrow$ 784 intensidades $\rightarrow$ una suma con pesos aprendidos $\rightarrow$ una puntuación entre 0 y 1 $\rightarrow$ una predicción $\rightarrow$ comparación con la referencia.}

\begin{itemize}
\item \textbf{Entrada:} 28 filas $\times$ 28 columnas de píxeles. Cada píxel contiene un número de brillo.
\item \textbf{Respuesta que queremos aprender:} la etiqueta publicada, \texttt{\small 0 =\allowbreak{} normal} o \texttt{\small 1 =\allowbreak{} neumonía}. La palabra «normal» es una categoría del conjunto.
\item \textbf{Aprender:} ajustar números llamados \textbf{pesos} usando imágenes y etiquetas de entrenamiento.
\item \textbf{Predecir:} aplicar esos pesos a otra imagen. La etiqueta de esa imagen se usa después para comprobar el resultado, no como entrada de la predicción.
\end{itemize}

Ejecuta de arriba abajo con \textbf{Mayús + Enter}, en una sesión de CPU. Las figuras guardadas permiten leer sin ejecutar; los controles aparecen al ejecutar su celda. Mueve un deslizador y suéltalo para recalcular su gráfico. Si editas \texttt{\small value=\allowbreak{}.\allowbreak{}.\allowbreak{}.\allowbreak{}} en el código, ejecuta de nuevo esa celda. Tras reiniciar la sesión, comienza desde el inicio.


\Sub{Tres conjuntos, tres trabajos diferentes}

\begin{tabularx}{\linewidth}{@{}YYY@{}}\toprule
Conjunto & Qué hacemos con él & Qué evitamos\\\midrule
Entrenamiento (\emph{train}) & El modelo aprende los pesos al comparar sus respuestas con las etiquetas. & Darle las imágenes reservadas para la evaluación final.\\
Validación (\emph{val}) & Comparamos tres ajustes del modelo y exploramos decisiones. & Elegir por el resultado de prueba.\\
Prueba (\emph{test}) & Medimos cómo responde el modelo ya elegido. & Volver a elegir modelo o umbral para mejorar esa cifra.\\\bottomrule\end{tabularx}

El informe usará un \textbf{umbral de 0.5}, fijado antes de evaluar prueba. Los controles sirven para explorar; no modifican esa decisión final. Es una práctica educativa con datos públicos de baja resolución. No estima riesgo clínico ni permite auditar independencia entre pacientes, porque el archivo no aporta sus identificadores.


\Pagina{BLOQUE 1}{Imágenes y herramientas}{}\label{bloque-1}

Primero abrimos los datos. Una \textbf{muestra} es una imagen; su \textbf{etiqueta de referencia} es la respuesta publicada que acompaña a esa imagen. Todavía no hay modelo entrenado. El archivo NPZ reúne varias matrices NumPy en un solo paquete.


\needspace{10\baselineskip}\phantomsection\addcontentsline{toc}{subsection}{1A · Cargar las herramientas de preparación}\Sub{1A · Cargar las herramientas de preparación}\label{paso-1a}

Cargamos las herramientas para abrir archivos y comprobar las bibliotecas. \texttt{\small import} hace que Python reconozca sus nombres; todavía no descarga datos ni entrena nada.

\texttt{\small Path} construye rutas, \texttt{\small urlopen} permite descargar bytes, \texttt{\small hashlib} calcula una huella del archivo y \texttt{\small json} escribirá el informe. \texttt{\small sys.\allowbreak{}executable} identifica el Python de esta sesión, \texttt{\small subprocess} ejecuta su instalador y \texttt{\small version} consulta versiones. La salida confirma que la preparación puede continuar.


\begin{Codigo}{1A · Código}
import sys, subprocess, importlib.util, hashlib, json
from pathlib import Path
from importlib.metadata import version
from urllib.request import urlopen
print("Herramientas de preparación disponibles.")
\end{Codigo}

\begin{Salida}{1A · Resultado}
Herramientas de preparación disponibles.
\end{Salida}

\needspace{10\baselineskip}\phantomsection\addcontentsline{toc}{subsection}{1B · Preparar las bibliotecas ausentes}\Sub{1B · Preparar las bibliotecas ausentes}\label{paso-1b}

Instalamos sólo bibliotecas ausentes. \texttt{\small requisitos} relaciona cada módulo con su paquete; la lista con \texttt{\small for .\allowbreak{}.\allowbreak{}.\allowbreak{} if} conserva los que \texttt{\small find\_\allowbreak{}spec} no encuentra.

\texttt{\small subprocess.\allowbreak{}run} llama a \texttt{\small pip} usando el Python activo; \texttt{\small capture\_\allowbreak{}output=\allowbreak{}True} recoge los mensajes. \texttt{\small if proceso.\allowbreak{}returncode} detecta un fallo y \texttt{\small raise} lo muestra. Si no falta ninguna biblioteca, no instala nada. Una instalación necesita internet y puede tardar.


\begin{Codigo}{1B · Código}
requisitos = {"numpy": "numpy>=1.24,<3", "matplotlib": "matplotlib>=3.8,<4", "sklearn": "scikit-learn>=1.4,<2", "ipywidgets": "ipywidgets>=8,<9"}
faltantes = [p for m, p in requisitos.items() if importlib.util.find_spec(m) is None]
if faltantes:
    proceso = subprocess.run([sys.executable, "-m", "pip", "install", "-q", *faltantes], capture_output=True, text=True)
    if proceso.returncode:
        raise RuntimeError("No terminó la instalación. Revisa internet. " + proceso.stderr[-1000:])
print("Preparación terminada; bibliotecas solicitadas:", faltantes or "ninguna")
\end{Codigo}

\begin{Salida}{1B · Resultado}
Preparación terminada; bibliotecas solicitadas: ninguna
\end{Salida}

\needspace{10\baselineskip}\phantomsection\addcontentsline{toc}{subsection}{1C · Importar matrices, gráficos y modelos}\Sub{1C · Importar matrices, gráficos y modelos}\label{paso-1c}

Ahora damos nombres a las herramientas que usarán las celdas. \texttt{\small numpy as np} calcula con matrices y \texttt{\small matplotlib.\allowbreak{}pyplot as plt} dibuja. \texttt{\small as} sólo abrevia el nombre.

\texttt{\small DummyClassifier} crea la regla que responde siempre lo mismo. \texttt{\small LogisticRegression} será el modelo que aprende de píxeles y etiquetas. Las métricas comparan respuestas; \texttt{\small widgets} crea controles y \texttt{\small display} los muestra. En Colab se habilita su visor con \texttt{\small enable\_\allowbreak{}custom\_\allowbreak{}widget\_\allowbreak{}manager(\allowbreak{})\allowbreak{}}. \textbf{Importar una clase no la entrena}: sólo permite crearla después.


\begin{Codigo}{1C · Código}
import numpy as np
import matplotlib.pyplot as plt
from sklearn.dummy import DummyClassifier
from sklearn.linear_model import LogisticRegression
from sklearn.metrics import accuracy_score, balanced_accuracy_score, confusion_matrix
import ipywidgets as widgets
from IPython.display import display
if "google.colab" in sys.modules:
    from google.colab import output
    output.enable_custom_widget_manager()
print("Bibliotecas y controles disponibles.")
\end{Codigo}

\begin{Salida}{1C · Resultado}
Bibliotecas y controles disponibles.
\end{Salida}

\needspace{10\baselineskip}\phantomsection\addcontentsline{toc}{subsection}{1D · Localizar o descargar los datos públicos}\Sub{1D · Localizar o descargar los datos públicos}\label{paso-1d}

Buscamos \texttt{\small datos/\allowbreak{}pneumoniamnist.\allowbreak{}npz}. \texttt{\small exists(\allowbreak{})\allowbreak{}} evita descargar otra copia si ya está ahí. Si falta, \texttt{\small mkdir(\allowbreak{}.\allowbreak{}.\allowbreak{}.\allowbreak{},\allowbreak{} parents=\allowbreak{}True)\allowbreak{}} crea la carpeta y \texttt{\small urlopen} recibe el archivo público; \texttt{\small with} cierra la conexión al terminar.

Una huella SHA-256 es un identificador calculado a partir de todos los bytes. \texttt{\small sha256(\allowbreak{}.\allowbreak{}.\allowbreak{}.\allowbreak{})\allowbreak{}.\allowbreak{}hexdigest(\allowbreak{})\allowbreak{}} la calcula y \texttt{\small assert} exige la huella esperada antes de \texttt{\small write\_\allowbreak{}bytes}. Esto comprueba que usamos el mismo archivo; no interpreta las imágenes.


\begin{Codigo}{1D · Código}
ruta = Path("datos/pneumoniamnist.npz")
sha_esperada = "e1792d3f03751cb101e99f19a63b3c1941436c988665f47853417b05be250cd8"
if not ruta.exists():
    ruta.parent.mkdir(parents=True, exist_ok=True)
    url = "https://zenodo.org/records/10519652/files/pneumoniamnist.npz?download=1"
    with urlopen(url, timeout=90) as respuesta:
        contenido = respuesta.read()
    assert hashlib.sha256(contenido).hexdigest() == sha_esperada, "La descarga no coincide con la copia verificada."
    ruta.write_bytes(contenido)
print("Archivo disponible:", ruta)
\end{Codigo}

\begin{Salida}{1D · Resultado}
Archivo disponible: datos\pneumoniamnist.npz
\end{Salida}

\needspace{10\baselineskip}\phantomsection\addcontentsline{toc}{subsection}{1E · Verificar y abrir el archivo}\Sub{1E · Verificar y abrir el archivo}\label{paso-1e}

Verificamos también la copia que ya existía y abrimos el paquete. \texttt{\small read\_\allowbreak{}bytes(\allowbreak{})\allowbreak{}} lee el archivo; \texttt{\small np.\allowbreak{}load(\allowbreak{}.\allowbreak{}.\allowbreak{}.\allowbreak{},\allowbreak{} allow\_\allowbreak{}pickle=\allowbreak{}False)\allowbreak{}} abre sus matrices sin cargar objetos de Python.

\texttt{\small paquete.\allowbreak{}files} lista las claves del paquete. Esperamos imágenes y etiquetas para entrenamiento, validación y prueba. Abrir este contenedor no significa que ya estemos usando prueba para elegir el modelo.


\begin{Codigo}{1E · Código}
assert hashlib.sha256(ruta.read_bytes()).hexdigest() == sha_esperada, "Revisa el archivo local: su huella es diferente."
paquete = np.load(ruta, allow_pickle=False)
assert set(paquete.files) == {f"{s}_{k}" for s in ("train", "val", "test") for k in ("images", "labels")}
print("Huella verificada. Arreglos:", sorted(paquete.files))
\end{Codigo}

\begin{Salida}{1E · Resultado}
Huella verificada. Arreglos: ['test_images', 'test_labels', 'train_images', 'train_labels', 'val_images', 'val_labels']
\end{Salida}

\needspace{10\baselineskip}\phantomsection\addcontentsline{toc}{subsection}{1F · Extraer entrenamiento y validación}\Sub{1F · Extraer entrenamiento y validación}\label{paso-1f}

Extraemos sólo entrenamiento y validación. \texttt{\small paquete[\allowbreak{}"train\_\allowbreak{}images"]\allowbreak{}} entrega los píxeles; \texttt{\small paquete[\allowbreak{}"train\_\allowbreak{}labels"]\allowbreak{}} entrega sus etiquetas. \texttt{\small .\allowbreak{}ravel(\allowbreak{})\allowbreak{}} deja las etiquetas en una lista numérica de una dimensión.

Lee \texttt{\small (\allowbreak{}4708,\allowbreak{} 28,\allowbreak{} 28)\allowbreak{}} como \textbf{4708 imágenes, cada una de 28 filas y 28 columnas}. Cada posición de \texttt{\small y\_\allowbreak{}train\_\allowbreak{}original} corresponde a la imagen de la misma posición. \texttt{\small uint8} significa enteros entre 0 y 255; \texttt{\small unique} comprueba que sólo existan etiquetas 0 y 1. Los \texttt{\small assert} detienen la ejecución si los datos no coinciden con lo previsto.


\begin{Codigo}{1F · Código}
train_original = paquete["train_images"]
y_train_original = paquete["train_labels"].ravel()
val_original = paquete["val_images"]
y_val_original = paquete["val_labels"].ravel()
assert train_original.shape == (4708, 28, 28) and val_original.shape == (524, 28, 28)
assert train_original.dtype == np.uint8 and set(np.unique(y_train_original)) == {0, 1}
print(f"Entrenamiento: {len(train_original)} imágenes de 28 × 28 píxeles y {len(y_train_original)} etiquetas.")
print(f"Validación: {len(val_original)} imágenes de 28 × 28 píxeles y {len(y_val_original)} etiquetas.")
print("Referencia publicada: 0 = normal; 1 = neumonía.")
\end{Codigo}

\begin{Salida}{1F · Resultado}
Entrenamiento: 4708 imágenes de 28 × 28 píxeles y 4708 etiquetas.
Validación: 524 imágenes de 28 × 28 píxeles y 524 etiquetas.
Referencia publicada: 0 = normal; 1 = neumonía.
\end{Salida}

\needspace{10\baselineskip}\phantomsection\addcontentsline{toc}{subsection}{1G · Preparar etiquetas y ejemplos de la galería}\Sub{1G · Preparar etiquetas y ejemplos de la galería}\label{paso-1g}

Elegimos tres ejemplos de cada etiqueta para que la galería muestre ambas clases. \texttt{\small y\_\allowbreak{}train\_\allowbreak{}original =\allowbreak{}=\allowbreak{} c} pregunta qué etiquetas son iguales a c; \texttt{\small flatnonzero} obtiene sus posiciones y \texttt{\small [\allowbreak{}:\allowbreak{}3]\allowbreak{}} conserva las tres primeras. \texttt{\small concatenate} reúne los dos grupos.

\texttt{\small nombres} traduce 0 y 1 a sus rótulos; \texttt{\small colores} asigna colores de dibujo y \texttt{\small rcParams.\allowbreak{}update} fija un estilo legible. Estos nombres y colores no son características que aprenda el modelo.


\begin{Codigo}{1G · Código}
nombres = np.array(["Normal (etiqueta 0)", "Neumonía (etiqueta 1)"])
colores = ["#226b8a", "#c96932"]
plt.rcParams.update({"font.size": 11, "axes.spines.top": False, "axes.spines.right": False})
indices = np.concatenate([np.flatnonzero(y_train_original == c)[:3] for c in (0, 1)])
print("Índices de la galería:", indices.tolist())
\end{Codigo}

\begin{Salida}{1G · Resultado}
Índices de la galería: [5, 7, 15, 0, 1, 2]
\end{Salida}

\needspace{10\baselineskip}\phantomsection\addcontentsline{toc}{subsection}{1H · Mostrar las primeras radiografías}\Sub{1H · Mostrar las primeras radiografías}\label{paso-1h}

Cada panel muestra una matriz de 28 $\times$ 28 valores y su \textbf{referencia publicada}, no una predicción. El aspecto borroso refleja la baja resolución del conjunto.

\texttt{\small subplots(\allowbreak{}2,\allowbreak{} 3)\allowbreak{}} crea seis paneles. \texttt{\small .\allowbreak{}ravel(\allowbreak{})\allowbreak{}} los recorre en una fila y \texttt{\small zip} empareja cada panel con una imagen. \texttt{\small imshow} convierte valores en grises; \texttt{\small vmin=\allowbreak{}0,\allowbreak{} vmax=\allowbreak{}255} mantiene la misma escala e \texttt{\small interpolation=\allowbreak{}"nearest"} evita inventar píxeles intermedios. \texttt{\small set\_\allowbreak{}title} escribe el rótulo, \texttt{\small axis(\allowbreak{}"off")\allowbreak{}} oculta ejes y \texttt{\small show} presenta la figura.


\begin{Codigo}{1H · Código}
fig, ejes = plt.subplots(2, 3, figsize=(10, 6.2), layout="constrained")
for ax, i in zip(ejes.ravel(), indices):
    ax.imshow(train_original[i], cmap="gray", vmin=0, vmax=255, interpolation="nearest")
    ax.set_title(f"{nombres[y_train_original[i]]}\nTrain original · índice {i}", color=colores[y_train_original[i]])
    ax.axis("off")
plt.show()
print("PneumoniaMNIST · archivo público verificado · 4 170 669 bytes · CPU")
print("Versiones: " + ", ".join(f"{p} {version(p)}" for p in ("numpy", "matplotlib", "scikit-learn")))
\end{Codigo}

\Figura{paso-1h-1.png}{1H · Seis imágenes de entrenamiento y sus etiquetas publicadas. Una etiqueta de referencia y los píxeles son datos distintos.}{94mm}

\begin{Salida}{1H · Resultado}
PneumoniaMNIST · archivo público verificado · 4 170 669 bytes · CPU
Versiones: numpy 2.5.4, matplotlib 3.11.2, scikit-learn 1.9.1
\end{Salida}

\Pagina{BLOQUE 2}{Entrenamiento, validación y prueba}{}\label{bloque-2}

Antes de aprender, quitamos copias idénticas: si la misma imagen aparece al aprender y al evaluar, el resultado puede parecer mejor por haberla visto. Conservamos juntos los píxeles y su etiqueta. El conjunto de prueba queda reservado hasta fijar las decisiones.


\needspace{10\baselineskip}\phantomsection\addcontentsline{toc}{subsection}{2A · Conservar la primera copia de entrenamiento}\Sub{2A · Conservar la primera copia de entrenamiento}\label{paso-2a}

Quitamos imágenes exactamente repetidas de entrenamiento y conservamos su primera aparición. \texttt{\small np.\allowbreak{}unique(\allowbreak{}.\allowbreak{}.\allowbreak{}.\allowbreak{},\allowbreak{} axis=\allowbreak{}0,\allowbreak{} return\_\allowbreak{}index=\allowbreak{}True)\allowbreak{}} compara las imágenes completas y devuelve sus primeras posiciones. \texttt{\small \_\allowbreak{}} recibe un resultado que no utilizaremos.

\texttt{\small np.\allowbreak{}sort} mantiene el orden original. Los mismos índices se aplican a imágenes y etiquetas para no mezclarlas. La salida debe indicar \textbf{4690 imágenes}, frente a 4708 originales: se excluyen 18 copias.


\begin{Codigo}{2A · Código}
_, primeras = np.unique(train_original, axis=0, return_index=True)
ids_train = np.sort(primeras)
imagenes_train = train_original[ids_train]
y_train = y_train_original[ids_train]
print("Entrenamiento conservado:", len(y_train), "· retiradas:", len(train_original)-len(y_train))
\end{Codigo}

\begin{Salida}{2A · Resultado}
Entrenamiento conservado: 4690 · retiradas: 18
\end{Salida}

\needspace{10\baselineskip}\phantomsection\addcontentsline{toc}{subsection}{2B · Excluir coincidencias de validación}\Sub{2B · Excluir coincidencias de validación}\label{paso-2b}

Comprobamos si una imagen de validación ya está en entrenamiento. \texttt{\small tobytes(\allowbreak{})\allowbreak{}} representa sus píxeles como bytes; el conjunto \texttt{\small \{.\allowbreak{}.\allowbreak{}.\allowbreak{} for .\allowbreak{}.\allowbreak{}.\allowbreak{}\}} permite buscar coincidencias exactas con \texttt{\small not in}.

\texttt{\small conservar\_\allowbreak{}val} es una lista de respuestas verdadero/falso. Al usarla entre corchetes seleccionamos las imágenes y etiquetas correspondientes. \texttt{\small \textasciitilde{}} invierte esas respuestas para contar exclusiones. Deben quedar \textbf{516 imágenes de validación}: se retiran 8 de las 524 originales.


\begin{Codigo}{2B · Código}
huellas_train = {imagen.tobytes() for imagen in imagenes_train}
conservar_val = np.array([imagen.tobytes() not in huellas_train for imagen in val_original])
ids_val = np.flatnonzero(conservar_val)
imagenes_val = val_original[conservar_val]
y_val = y_val_original[conservar_val]
print("Validación conservada:", len(y_val), "· excluidas:", int((~conservar_val).sum()))
\end{Codigo}

\begin{Salida}{2B · Resultado}
Validación conservada: 516 · excluidas: 8
\end{Salida}

\needspace{10\baselineskip}\phantomsection\addcontentsline{toc}{subsection}{2C · Comprobar cantidades y etiquetas}\Sub{2C · Comprobar cantidades y etiquetas}\label{paso-2c}

Confirmamos cuántas imágenes quedan y cuántas pertenecen a cada etiqueta. \texttt{\small bincount} cuenta los ceros y los unos; \texttt{\small array\_\allowbreak{}equal} exige los conteos esperados. Un \texttt{\small assert} sin error significa que la condición se cumplió.

No confundas esta auditoría con separar pacientes: dos imágenes distintas de una misma persona no son copias exactas. Aquí no tenemos identificadores para comprobar esa independencia.


\begin{Codigo}{2C · Código}
assert len(y_train) == 4690 and len(y_val) == 516
assert np.array_equal(np.bincount(y_train), [1213, 3477])
assert np.array_equal(np.bincount(y_val), [135, 381])
print("Conteos por etiqueta: train", np.bincount(y_train), "· val", np.bincount(y_val))
\end{Codigo}

\begin{Salida}{2C · Resultado}
Conteos por etiqueta: train [1213 3477] · val [135 381]
\end{Salida}

\needspace{10\baselineskip}\phantomsection\addcontentsline{toc}{subsection}{2D · Visualizar las particiones}\Sub{2D · Visualizar las particiones}\label{paso-2d}

Las barras muestran que las clases de entrenamiento tienen tamaños diferentes. Las tarjetas explican el trabajo de cada grupo; \textbf{sólo entrenamiento ajusta pesos}.

\texttt{\small bar} dibuja los conteos y \texttt{\small bar\_\allowbreak{}label} los escribe. \texttt{\small set} rotula los ejes. En el segundo panel, \texttt{\small text} escribe las tarjetas; \texttt{\small transAxes} sitúa cada texto con coordenadas de 0 a 1 y \texttt{\small bbox} dibuja su fondo. Los mensajes inferiores indican cuántas imágenes se retiraron y qué índices originales fueron excluidos de validación.


\begin{Codigo}{2D · Código}
fig, ejes = plt.subplots(1, 2, figsize=(11, 4.3), layout="constrained")
conteos = np.bincount(y_train, minlength=2)
barras = ejes[0].bar(["Normal", "Neumonía"], conteos, color=colores)
ejes[0].bar_label(barras, padding=3)
ejes[0].set(title="Entrenamiento tras quitar copias", ylabel="Número de imágenes", ylim=(0, 4000))
ejes[1].axis("off")
for y, texto, color in [(0.82, "ENTRENAR · 4690\nAprender los pesos", "#226b8a"), (0.48, "VALIDAR · 516\nElegir C entre tres opciones", "#007f89"), (0.14, "PROBAR · 624\nEvaluación final con decisiones fijadas", "#c96932")]:
    ejes[1].text(0.5, y, texto, ha="center", va="center", color="white", transform=ejes[1].transAxes, bbox=dict(boxstyle="round,pad=0.7", facecolor=color, edgecolor="none"))
plt.show()
print("Eliminadas de train:", len(train_original)-len(y_train), "· excluidas de val:", int((~conservar_val).sum()))
print("Índices originales de val excluidos:", np.flatnonzero(~conservar_val).tolist())
\end{Codigo}

\Figura{paso-2d-1.png}{2D · Entrenamiento tiene 4690 imágenes sin copias exactas. Validación conserva 516 y prueba reserva 624; cada grupo cumple una función distinta.}{94mm}

\begin{Salida}{2D · Resultado}
Eliminadas de train: 18 · excluidas de val: 8
Índices originales de val excluidos: [61, 105, 172, 236, 244, 253, 302, 473]
\end{Salida}

\Pagina{BLOQUE 3}{De imagen a características}{}\label{bloque-3}

El modelo recibe una tabla de números llamada \textbf{X}. Cada fila corresponde a una imagen y cada columna a una intensidad de píxel: 28 $\times$ 28 = \textbf{784 características}. La \textbf{y} minúscula guarda las etiquetas, una por fila. Aplanar reorganiza los valores; no añade información ni incorpora la etiqueta a X.


\needspace{10\baselineskip}\phantomsection\addcontentsline{toc}{subsection}{3A · Preparar las características de entrenamiento}\Sub{3A · Preparar las características de entrenamiento}\label{paso-3a}

Convertimos cada imagen en una fila de 784 números. \texttt{\small reshape(\allowbreak{}len(\allowbreak{}imagenes\_\allowbreak{}train)\allowbreak{},\allowbreak{} -1)\allowbreak{}} conserva el número de imágenes; \texttt{\small -1} pide a NumPy calcular cuántas columnas faltan. Recorre primero una fila de píxeles, después la siguiente.

\texttt{\small astype(\allowbreak{}np.\allowbreak{}float64)\allowbreak{}} permite decimales y \texttt{\small /\allowbreak{} 255.\allowbreak{}0} lleva el brillo al intervalo 0–1: 0 pasa a 0; 255 pasa a 1; 128 pasa a aproximadamente 0.502. La salida \texttt{\small (\allowbreak{}4690,\allowbreak{} 784)\allowbreak{}} significa \textbf{4690 ejemplos con 784 entradas cada uno}. \texttt{\small y\_\allowbreak{}train} permanece aparte.


\begin{Codigo}{3A · Código}
X_train = imagenes_train.reshape(len(imagenes_train), -1).astype(np.float64) / 255.0
print("X_train:", X_train.shape, "· valores por imagen:", X_train.shape[1])
\end{Codigo}

\begin{Salida}{3A · Resultado}
X_train: (4690, 784) · valores por imagen: 784
\end{Salida}

\needspace{10\baselineskip}\phantomsection\addcontentsline{toc}{subsection}{3B · Aplicar la misma regla a validación}\Sub{3B · Aplicar la misma regla a validación}\label{paso-3b}

Transformamos validación con la misma regla: mismo orden de píxeles y mismo divisor 255. Si cambiáramos el orden, una columna representaría otro lugar de la imagen y el modelo recibiría algo distinto de lo que aprendió.

Los \texttt{\small assert} comprueban \texttt{\small (\allowbreak{}4690,\allowbreak{} 784)\allowbreak{}} para entrenamiento y \texttt{\small (\allowbreak{}516,\allowbreak{} 784)\allowbreak{}} para validación. No estamos calculando estadísticas nuevas con validación; sólo reorganizamos y dividimos sus valores.


\begin{Codigo}{3B · Código}
X_val = imagenes_val.reshape(len(imagenes_val), -1).astype(np.float64) / 255.0
assert X_train.shape == (4690, 784) and X_val.shape == (516, 784)
print("X_val:", X_val.shape, "· transformaciones comprobadas")
\end{Codigo}

\begin{Salida}{3B · Resultado}
X_val: (516, 784) · transformaciones comprobadas
\end{Salida}

\needspace{10\baselineskip}\phantomsection\addcontentsline{toc}{subsection}{3C · Elegir una imagen para inspeccionar}\Sub{3C · Elegir una imagen para inspeccionar}\label{paso-3c}

\texttt{\small i =\allowbreak{} 0} selecciona el primer ejemplo; Python empieza a contar en cero. \texttt{\small X\_\allowbreak{}train[\allowbreak{}i]\allowbreak{}} es su fila de 784 intensidades y \texttt{\small y\_\allowbreak{}train[\allowbreak{}i]\allowbreak{}} es su referencia.

El \texttt{\small assert} evita un índice fuera de rango. En el visor podrás cambiar \textbf{Imagen} para verificar que los píxeles, el vector y la etiqueta cambian juntos.


\begin{Codigo}{3C · Código}
i = 0
assert 0 <= i < len(X_train), "Elige un índice de entrenamiento válido."
print("Imagen elegida:", i, "· etiqueta:", y_train[i])
\end{Codigo}

\begin{Salida}{3C · Resultado}
Imagen elegida: 0 · etiqueta: 1
\end{Salida}

\needspace{10\baselineskip}\phantomsection\addcontentsline{toc}{subsection}{3D · Relacionar imagen, vector y etiqueta}\Sub{3D · Relacionar imagen, vector y etiqueta}\label{paso-3d}

\textbf{Cambia Imagen.} A la izquierda ves la radiografía, en el centro sus 784 intensidades ordenadas y a la derecha su etiqueta separada. La curva no es una señal temporal: su eje horizontal es la posición del píxel al aplanar.

El mensaje muestra una comprobación real: el píxel de fila 10, columna 10 termina en la posición \texttt{\small 10*28 + 10 =\allowbreak{} 290} del vector, dividido entre 255. No se pierde su valor por reorganizarlo.

\texttt{\small def ver\_\allowbreak{}imagen} define qué redibujar; \texttt{\small imshow} pinta la imagen y \texttt{\small plot} el vector. \texttt{\small scatter} marca el mismo píxel en ambas vistas con color naranja. \texttt{\small IntSlider} fija rango y valor inicial. \texttt{\small interactive} conecta el control con la función; \texttt{\small continuous\_\allowbreak{}update=\allowbreak{}False} recalcula al soltarlo y \texttt{\small display} muestra el visor.


\begin{Codigo}{3D · Código}
def ver_imagen(indice_imagen=0):
    i = indice_imagen
    fig, ejes = plt.subplots(1, 3, figsize=(12, 3.8), layout="constrained", gridspec_kw={"width_ratios": [1, 1.8, 1.4]})
    ejes[0].imshow(imagenes_train[i], cmap="gray", vmin=0, vmax=255, interpolation="nearest")
    ejes[0].scatter(10, 10, s=85, facecolors="none", edgecolors="#c96932", linewidths=2)
    ejes[0].set(title="Una muestra: 28 × 28", xlabel="Columna", ylabel="Fila")
    ejes[1].plot(np.arange(784), X_train[i], color="#007f89", linewidth=0.9)
    ejes[1].scatter(290, X_train[i, 290], s=70, color="#c96932", zorder=3)
    ejes[1].set(title="Una fila de X: 784 valores", xlabel="Posición del píxel al aplanar", ylabel="Intensidad / 255", ylim=(0, 1))
    ejes[2].axis("off")
    ejes[2].text(0.05, 0.75, f"X_train: {X_train.shape}\n\n4690 filas = imágenes\n784 columnas = píxeles\n\ny_train[{i}] = {y_train[i]}\n{nombres[y_train[i]]}", va="top", transform=ejes[2].transAxes)
    plt.show()
    fila_pixel, columna_pixel = 10, 10
    posicion = fila_pixel * 28 + columna_pixel
    brillo = int(imagenes_train[i, fila_pixel, columna_pixel])
    print(f"Imagen {i}: referencia {y_train[i]} = {nombres[y_train[i]]}.")
    print(f"Píxel [10, 10] = {brillo} → X_train[{i}, 290] = {brillo}/255 = {X_train[i, posicion]:.6f}")
    assert np.isclose(X_train[i, posicion], brillo / 255.0)
    return fig

imagen_control = widgets.IntSlider(value=i, min=0, max=len(X_train)-1, step=1,
                                   description="Imagen", continuous_update=False)
visor_imagen = widgets.interactive(ver_imagen, indice_imagen=imagen_control)
display(visor_imagen)
\end{Codigo}

\Figura{paso-3d-1.png}{3D · El píxel [10,10] aparece en la posición 290 del vector. Dividir su intensidad entre 255 cambia la escala, no su posición ni su etiqueta.}{94mm}

\begin{Salida}{3D · Resultado}
Imagen 0: referencia 1 = Neumonía (etiqueta 1).
Píxel [10, 10] = 164 → X_train[0, 290] = 164/255 = 0.643137
\end{Salida}

\Pagina{BLOQUE 4}{Cómo una suma se convierte en puntuación}{}\label{bloque-4}

La regresión logística multiplica cada característica por un \textbf{peso}, suma los productos y añade un \textbf{sesgo}: \texttt{\small z =\allowbreak{} x1*w1 + x2*w2 + .\allowbreak{}.\allowbreak{}.\allowbreak{} + b}. Los pesos y b son números que el entrenamiento aprenderá. Una función llamada \textbf{sigmoide} transforma z en una puntuación entre 0 y 1. Aquí practicamos esa cuenta antes de entrenar.


\Sub{Una cuenta con cuatro entradas ficticias}\Figura{suma_cuatro_entradas.png}{Entradas y pesos inventados: cada posición aporta un producto. No son valores de radiografías ni pesos aprendidos con pacientes.}{82mm}
\begin{enumerate}
\item \textbf{Multiplica cada entrada por su peso.} La segunda contribución es negativa: \(0,80\times(-1)=-0,80\). Un peso puede ser negativo o mayor que uno; no es una intensidad de píxel.
\item \textbf{Suma productos y añade el término independiente.} Los productos suman 0,80. Añadir \(b=-0,30\) deja \(z=0,50\). El término independiente desplaza la suma; no multiplica una entrada.
\item \textbf{Aplica la sigmoide y el umbral.} La fórmula transforma \(z=0,50\) en una puntuación de aproximadamente 0,622. Como \(0,622\geq0,5\), la predicción es clase 1.
\end{enumerate}
\begin{Codigo}{Cuenta ilustrativa · no modifica el modelo del taller}
x_ejemplo = np.array([0.20, 0.80, 0.60, 0.40])
w_ejemplo = np.array([1.00, -1.00, 2.00, 0.50])
productos_ejemplo = x_ejemplo * w_ejemplo
z_ejemplo = productos_ejemplo.sum() - 0.30
s_ejemplo = 1 / (1 + np.exp(-z_ejemplo))
# productos: [0.20, -0.80, 1.20, 0.20]; z: 0.50; s: 0.622
\end{Codigo}
\texttt{*} multiplica las posiciones correspondientes; \texttt{.sum()} reúne sus aportaciones. \textbf{Predecir} reutiliza pesos; \textbf{aprender} ajusta pesos a partir de ejemplos. Aquí los elegimos a mano. Un modelo entrenado ajusta sus pesos con imágenes y etiquetas.
\begin{Practica}{Cambia una entrada, conserva los pesos}
Si la tercera entrada baja de 0,60 a 0,10, su contribución baja de 1,20 a 0,20. Queda \(z=-0,50\), la puntuación es aproximadamente 0,378 y la predicción pasa a clase 0. No se ha reentrenado. Es un ejemplo matemático, no una interpretación anatómica.
\end{Practica}



\needspace{10\baselineskip}\phantomsection\addcontentsline{toc}{subsection}{4A · Elegir sumas para explorar}\Sub{4A · Elegir sumas para explorar}\label{paso-4a}

Practiquemos una suma pequeña con \textbf{dos entradas ficticias}, no con una radiografía: \texttt{\small x=\allowbreak{}[\allowbreak{}0.\allowbreak{}25,\allowbreak{} 0.\allowbreak{}75]\allowbreak{}}, pesos \texttt{\small [\allowbreak{}1,\allowbreak{} -1]\allowbreak{}} y sesgo \texttt{\small b=\allowbreak{}0.\allowbreak{}5}. \texttt{\small np.\allowbreak{}dot} multiplica valores de la misma posición y suma: \texttt{\small 0.\allowbreak{}25*1 + 0.\allowbreak{}75*(\allowbreak{}-1)\allowbreak{} + 0.\allowbreak{}5 =\allowbreak{} 0}.

Estos pesos son elegidos para explicar la cuenta. El modelo real aprenderá 784 pesos. \texttt{\small linspace(\allowbreak{}-6,\allowbreak{} 6,\allowbreak{} 400)\allowbreak{}} prepara muchas sumas posibles para dibujar cómo se transforman.


\begin{Codigo}{4A · Código}
x_demo = np.array([0.25, 0.75])
w_demo = np.array([1.0, -1.0])
b_demo = 0.5
z_demo = float(np.dot(x_demo, w_demo) + b_demo)
z = np.linspace(-6, 6, 400)
print(f"Ejemplo ficticio: 0.25×1 + 0.75×(-1) + 0.5 = {z_demo:.1f}")
print(f"Para dibujar: {z.size} sumas posibles desde {z[0]} hasta {z[-1]}.")
\end{Codigo}

\begin{Salida}{4A · Resultado}
Ejemplo ficticio: 0.25×1 + 0.75×(-1) + 0.5 = 0.0
Para dibujar: 400 sumas posibles desde -6.0 hasta 6.0.
\end{Salida}

\needspace{10\baselineskip}\phantomsection\addcontentsline{toc}{subsection}{4B · Transformar las sumas con la sigmoide}\Sub{4B · Transformar las sumas con la sigmoide}\label{paso-4b}

La \textbf{sigmoide} calcula \texttt{\small p =\allowbreak{} 1 /\allowbreak{} (\allowbreak{}1 + exp(\allowbreak{}-z)\allowbreak{})\allowbreak{}}. \texttt{\small np.\allowbreak{}exp} es la función exponencial aplicada a cada suma. Para z = 0, \texttt{\small exp(\allowbreak{}0)\allowbreak{}=\allowbreak{}1}, por lo tanto \texttt{\small p=\allowbreak{}1/\allowbreak{}(\allowbreak{}1+1)\allowbreak{}=\allowbreak{}0.\allowbreak{}5}.

Una suma negativa produce p menor que 0.5 y una positiva, mayor que 0.5. La salida guarda una puntuación por cada z. En este clasificador p representará la salida para \textbf{etiqueta 1}; aún no es una etiqueta 0/1 ni una certeza clínica.


\begin{Codigo}{4B · Código}
p = 1 / (1 + np.exp(-z))
print("Rango de puntuaciones:", round(float(p.min()), 6), "a", round(float(p.max()), 6))
\end{Codigo}

\begin{Salida}{4B · Resultado}
Rango de puntuaciones: 0.002473 a 0.997527
\end{Salida}

\needspace{10\baselineskip}\phantomsection\addcontentsline{toc}{subsection}{4C · Ver la curva y sus ejemplos}\Sub{4C · Ver la curva y sus ejemplos}\label{paso-4c}

\textbf{Mueve Suma z} a -2, 0 y 2: las puntuaciones son aproximadamente 0.119, 0.500 y 0.881. El control sólo cambia una suma de ejemplo; todavía no aprende pesos.

\texttt{\small plot(\allowbreak{}z,\allowbreak{} p)\allowbreak{}} dibuja la curva y \texttt{\small exp(\allowbreak{}-suma\_\allowbreak{}z)\allowbreak{}} calcula el punto elegido. \texttt{\small scatter} lo marca; \texttt{\small annotate} muestra su cuenta y \texttt{\small axhline(\allowbreak{}0.\allowbreak{}5)\allowbreak{}} señala la mitad de la escala. \texttt{\small FloatSlider} admite decimales; \texttt{\small interactive} llama otra vez a \texttt{\small ver\_\allowbreak{}sigmoide}. Lee la secuencia impresa: \textbf{suma $\rightarrow$ puntuación $\rightarrow$ decisión si usamos umbral 0.5}.


\begin{Codigo}{4C · Código}
def ver_sigmoide(suma_z=0.0):
    fig, ax = plt.subplots(figsize=(9, 4.4), layout="constrained")
    ax.plot(z, p, color="#007f89", linewidth=3)
    punto = 1 / (1 + np.exp(-suma_z))
    ax.scatter(suma_z, punto, color="#c96932", s=85, zorder=3)
    ax.annotate(f"z={suma_z:.1f} → {punto:.3f}", (suma_z, punto), xytext=(8, 14), textcoords="offset points")
    ax.axhline(0.5, color="#697887", linestyle="--")
    ax.set(title="Sigmoide: mueve la suma y observa su puntuación", xlabel="Suma z", ylabel="Puntuación", xlim=(-6, 6), ylim=(-0.03, 1.12))
    plt.show()
    print(f"Suma z = {suma_z:.1f} → sigmoide = 1 / (1 + exp(-({suma_z:.1f}))) = {punto:.6f}")
    print(f"Con umbral 0.5: {punto:.6f} >= 0.5 es {punto >= 0.5} → predicción {int(punto >= 0.5)}.")
    return fig

suma_control = widgets.FloatSlider(value=0.0, min=-6.0, max=6.0, step=0.1,
                                  description="Suma z", continuous_update=False)
visor_sigmoide = widgets.interactive(ver_sigmoide, suma_z=suma_control)
display(visor_sigmoide)
\end{Codigo}

\Figura{paso-4c-1.png}{4C · La sigmoide transforma una suma z en una puntuación entre 0 y 1. Para z=0, la puntuación es 0.5.}{94mm}

\begin{Salida}{4C · Resultado}
Suma z = 0.0 → sigmoide = 1 / (1 + exp(-(0.0))) = 0.500000
Con umbral 0.5: 0.500000 >= 0.5 es True → predicción 1.
\end{Salida}

\Pagina{BLOQUE 5}{Una regla de referencia}{}\label{bloque-5}

¿Qué consigue un programa que responde siempre lo mismo? Esa será nuestra comparación mínima. Hay más imágenes etiquetadas como neumonía que como normal; por eso acertar muchas veces no significa reconocer bien ambas clases.


\needspace{10\baselineskip}\phantomsection\addcontentsline{toc}{subsection}{5A · Ajustar la regla mayoritaria}\Sub{5A · Ajustar la regla mayoritaria}\label{paso-5a}

Creamos una comparación que siempre responde con la clase más abundante. \texttt{\small DummyClassifier(\allowbreak{}strategy=\allowbreak{}"most\_\allowbreak{}frequent")\allowbreak{}} define esa regla y \texttt{\small .\allowbreak{}fit(\allowbreak{}X\_\allowbreak{}train,\allowbreak{} y\_\allowbreak{}train)\allowbreak{}} averigua cuál es la etiqueta mayoritaria.

Aquí \texttt{\small fit} no aprende relaciones entre píxeles y etiquetas. La regla obtiene su respuesta contando las etiquetas. Sirve para saber qué lograríamos sin analizar ninguna imagen.


\begin{Codigo}{5A · Código}
referencia = DummyClassifier(strategy="most_frequent")
referencia.fit(X_train, y_train)
print("Referencia ajustada únicamente con entrenamiento.")
\end{Codigo}

\begin{Salida}{5A · Resultado}
Referencia ajustada únicamente con entrenamiento.
\end{Salida}

\needspace{10\baselineskip}\phantomsection\addcontentsline{toc}{subsection}{5B · Predecir con la referencia}\Sub{5B · Predecir con la referencia}\label{paso-5b}

\texttt{\small referencia.\allowbreak{}predict(\allowbreak{}X\_\allowbreak{}val)\allowbreak{}} repite esa respuesta para cada imagen de validación. No reajusta la regla ni consulta \texttt{\small y\_\allowbreak{}val} al predecir.

\texttt{\small np.\allowbreak{}unique(\allowbreak{}pred\_\allowbreak{}base\_\allowbreak{}val)\allowbreak{}} muestra que todas las respuestas pertenecen a una sola clase. Ahora sí usaremos las referencias de validación para contar cuántas de esas respuestas coinciden.


\begin{Codigo}{5B · Código}
pred_base_val = referencia.predict(X_val)
print("Predicciones:", len(pred_base_val), "· clases predichas:", np.unique(pred_base_val))
\end{Codigo}

\begin{Salida}{5B · Resultado}
Predicciones: 516 · clases predichas: [1]
\end{Salida}

\needspace{10\baselineskip}\phantomsection\addcontentsline{toc}{subsection}{5C · Medir la referencia}\Sub{5C · Medir la referencia}\label{paso-5c}

\texttt{\small accuracy\_\allowbreak{}score(\allowbreak{}y\_\allowbreak{}val,\allowbreak{} pred\_\allowbreak{}base\_\allowbreak{}val)\allowbreak{}} calcula \textbf{coincidencias / todas las imágenes}. La regla mayoritaria acierta 381 de 516: aproximadamente 73.84\%, porque 381 referencias son de clase 1.

\texttt{\small balanced\_\allowbreak{}accuracy\_\allowbreak{}score} calcula el porcentaje reconocido en cada clase y promedia ambos. Esta regla reconoce 0\% de la clase 0 y 100\% de la clase 1: \texttt{\small (\allowbreak{}0 + 100)\allowbreak{}/\allowbreak{}2 =\allowbreak{} 50\%}. Así se ve por qué una sola cifra de exactitud puede ocultar que falla una clase completa. \texttt{\small .\allowbreak{}2\%} presenta una proporción como porcentaje.


\begin{Codigo}{5C · Código}
base_val = [accuracy_score(y_val, pred_base_val), balanced_accuracy_score(y_val, pred_base_val)]
print(f"Exactitud {base_val[0]:.2%} · equilibrada {base_val[1]:.2%}")
\end{Codigo}

\begin{Salida}{5C · Resultado}
Exactitud 73.84% · equilibrada 50.00%
\end{Salida}

\needspace{10\baselineskip}\phantomsection\addcontentsline{toc}{subsection}{5D · Comparar las dos medidas}\Sub{5D · Comparar las dos medidas}\label{paso-5d}

Las dos barras describen \textbf{el mismo conjunto de predicciones}. La primera cuenta imágenes; la segunda da el mismo peso a las dos clases. No están midiendo dos modelos diferentes.

\texttt{\small np.\allowbreak{}array(\allowbreak{}base\_\allowbreak{}val)\allowbreak{}*100} convierte proporciones en porcentajes; \texttt{\small bar} dibuja y \texttt{\small bar\_\allowbreak{}label} escribe las cifras. La salida añade la cuenta por clase. Usaremos exactitud equilibrada para comparar modelos porque interesa reconocer ambas etiquetas.


\begin{Codigo}{5D · Código}
fig, ax = plt.subplots(figsize=(8.5, 4), layout="constrained")
barras = ax.bar(["Exactitud", "Exactitud equilibrada"], np.array(base_val)*100, color=["#697887", "#c96932"])
ax.bar_label(barras, labels=[f"{v:.1%}" for v in base_val], padding=4)
ax.set(title="Validación · siempre predecir la clase mayoritaria", ylabel="Porcentaje", ylim=(0, 105))
plt.show()
print("Predicción constante:", nombres[int(pred_base_val[0])])
print(f"Exactitud = {int((pred_base_val == y_val).sum())}/{len(y_val)} = {base_val[0]:.2%}.")
print("Clase 0 reconocida: 0%. Clase 1 reconocida: 100%. Promedio: 50%.")
\end{Codigo}

\Figura{paso-5d-1.png}{5D · La regla mayoritaria acierta 73.84\% en validación, pero su exactitud equilibrada es 50\%: responde siempre con una clase.}{94mm}

\begin{Salida}{5D · Resultado}
Predicción constante: Neumonía (etiqueta 1)
Exactitud = 381/516 = 73.84%.
Clase 0 reconocida: 0%. Clase 1 reconocida: 100%. Promedio: 50%.
\end{Salida}

\Pagina{BLOQUE 6}{Entrenar y elegir C}{}\label{bloque-6}

Ahora sí aprendemos de las imágenes: \texttt{\small fit(\allowbreak{}X\_\allowbreak{}train,\allowbreak{} y\_\allowbreak{}train)\allowbreak{}} ajusta los 784 pesos y un sesgo. Probamos tres valores de \textbf{C}, que regulan cuánto se penalizan pesos grandes. Validación elige entre esos tres modelos ya entrenados. La etiqueta de validación sirve para medirlos, no para ajustar sus pesos.


\needspace{10\baselineskip}\phantomsection\addcontentsline{toc}{subsection}{6A · Fijar las opciones de regularización}\Sub{6A · Fijar las opciones de regularización}\label{paso-6a}

\texttt{\small opciones\_\allowbreak{}C =\allowbreak{} [\allowbreak{}0.\allowbreak{}01,\allowbreak{} 0.\allowbreak{}1,\allowbreak{} 1.\allowbreak{}0]\allowbreak{}} define tres candidatos antes de ver prueba. \textbf{Regularizar} significa penalizar pesos grandes durante el ajuste: limita cuánto puede depender la respuesta de combinaciones muy específicas del entrenamiento.

En esta implementación, menor C implica mayor penalización. C no es un porcentaje ni la puntuación de una imagen. No suponemos qué valor será mejor: lo comprobaremos con validación.


\begin{Codigo}{6A · Código}
opciones_C = [0.01, 0.1, 1.0]
print("Opciones de C:", opciones_C)
\end{Codigo}

\begin{Salida}{6A · Resultado}
Opciones de C: [0.01, 0.1, 1.0]
\end{Salida}

\needspace{10\baselineskip}\phantomsection\addcontentsline{toc}{subsection}{6B · Entrenar los tres candidatos}\Sub{6B · Entrenar los tres candidatos}\label{paso-6b}

El bucle \texttt{\small for} crea y entrena un modelo por cada C. \textbf{\texttt{\small candidato.\allowbreak{}fit(\allowbreak{}X\_\allowbreak{}train,\allowbreak{} y\_\allowbreak{}train)\allowbreak{}} es la operación que aprende los pesos} comparando imágenes y respuestas de entrenamiento.

\texttt{\small solver=\allowbreak{}"lbfgs"} elige el algoritmo que busca esos pesos; \texttt{\small max\_\allowbreak{}iter=\allowbreak{}2000} limita sus intentos y \texttt{\small tol=\allowbreak{}1e-4} fija cuándo puede detenerse. \texttt{\small class\_\allowbreak{}weight=\allowbreak{}"balanced"} da más peso al grupo menos frecuente durante el ajuste. \texttt{\small predict} aplica después cada modelo, sin reentrenarlo; las métricas cuentan su rendimiento. \texttt{\small append} conserva modelos y cifras. El \texttt{\small assert} avisa si se agotó el límite de iteraciones.


\begin{Codigo}{6B · Código}
modelos, resultados_C = [], []
for C in opciones_C:
    candidato = LogisticRegression(C=C, solver="lbfgs", max_iter=2000, class_weight="balanced", tol=1e-4)
    candidato.fit(X_train, y_train)
    assert candidato.n_iter_[0] < 2000, "El ajuste llegó al límite; revisa la convergencia antes de comparar."
    score_train = balanced_accuracy_score(y_train, candidato.predict(X_train))
    score_val = balanced_accuracy_score(y_val, candidato.predict(X_val))
    modelos.append(candidato)
    resultados_C.append([C, score_train, score_val])
    print(f"C={C:g}: pesos aprendidos; validación = {score_val:.2%} de exactitud equilibrada.")
print("Candidatos entrenados:", len(modelos), "· resultados nuevos:", len(resultados_C))
\end{Codigo}

\begin{Salida}{6B · Resultado}
C=0.01: pesos aprendidos; validación = 92.99% de exactitud equilibrada.
C=0.1: pesos aprendidos; validación = 95.05% de exactitud equilibrada.
C=1: pesos aprendidos; validación = 94.91% de exactitud equilibrada.
Candidatos entrenados: 3 · resultados nuevos: 3
\end{Salida}

\needspace{10\baselineskip}\phantomsection\addcontentsline{toc}{subsection}{6C · Organizar la tabla de resultados}\Sub{6C · Organizar la tabla de resultados}\label{paso-6c}

Cada fila corresponde a uno de los tres C. \texttt{\small np.\allowbreak{}array} convierte los resultados en una tabla numérica; el bucle de impresión añade encabezados y porcentajes para leerla.

\textbf{Entrenamiento} mide respuestas sobre ejemplos usados para aprender. \textbf{Validación} mide respuestas sobre otro grupo. Compara la columna de validación; una cifra de entrenamiento muy alta no basta para elegir. El código conserva las proporciones originales, sin redondearlas para la selección.


\begin{Codigo}{6C · Código}
resultados_C = np.array(resultados_C)
print("C       Entrenamiento    Validación   (exactitud equilibrada)")
for C, valor_train, valor_val in resultados_C:
    print(f"{C:<7g} {valor_train:>12.2%} {valor_val:>13.2%}")
\end{Codigo}

\begin{Salida}{6C · Resultado}
C       Entrenamiento    Validación   (exactitud equilibrada)
0.01          93.52%        92.99%
0.1           95.37%        95.05%
1             97.15%        94.91%
\end{Salida}

\needspace{10\baselineskip}\phantomsection\addcontentsline{toc}{subsection}{6D · Elegir mediante validación}\Sub{6D · Elegir mediante validación}\label{paso-6d}

\texttt{\small resultados\_\allowbreak{}C[\allowbreak{}:\allowbreak{},\allowbreak{} 2]\allowbreak{}} toma todas las filas de la tercera columna: validación. \texttt{\small argmax} encuentra la fila de mayor valor; ese índice selecciona el modelo y su C.

En la ejecución guardada gana \textbf{C = 0.1}. Si tus bibliotecas producen otra elección, registra tu resultado: no mires prueba para forzar una coincidencia. \texttt{\small modelo.\allowbreak{}coef\_\allowbreak{}} contiene los \textbf{784 pesos aprendidos} y \texttt{\small intercept\_\allowbreak{}}, el sesgo. Estos números se obtuvieron con entrenamiento; C se eligió con validación.


\begin{Codigo}{6D · Código}
mejor = int(np.argmax(resultados_C[:, 2]))
modelo = modelos[mejor]
C_elegido = float(resultados_C[mejor, 0])
print(f"Elegido: C={C_elegido:g} · validación {resultados_C[mejor, 2]:.2%}")
print(f"Modelo elegido: {modelo.coef_.shape[1]} pesos y un sesgo b = {modelo.intercept_[0]:.6f}.")
\end{Codigo}

\begin{Salida}{6D · Resultado}
Elegido: C=0.1 · validación 95.05%
Modelo elegido: 784 pesos y un sesgo b = 1.088078.
\end{Salida}

\needspace{10\baselineskip}\phantomsection\addcontentsline{toc}{subsection}{6E · Visualizar y registrar la comparación}\Sub{6E · Visualizar y registrar la comparación}\label{paso-6e}

\textbf{Pulsa uno de los valores de C} para señalar sus resultados. El marcador azul cambia de candidato; el círculo naranja sigue indicando el modelo elegido por validación. Inspeccionar no vuelve a entrenar ni cambia la elección del informe.

\texttt{\small plot} dibuja las dos columnas, \texttt{\small flatnonzero} localiza la fila seleccionada y \texttt{\small scatter} añade los marcadores. \texttt{\small xscale=\allowbreak{}"log"} separa 0.01, 0.1 y 1 por factores de diez. \texttt{\small ToggleButtons} presenta los tres valores como botones e \texttt{\small interactive} redibuja la vista.


\begin{Codigo}{6E · Código}
def ver_candidatos(C_vista=C_elegido):
    fila = int(np.flatnonzero(resultados_C[:, 0] == C_vista)[0])
    fig, ax = plt.subplots(figsize=(9, 4.2), layout="constrained")
    ax.plot(opciones_C, resultados_C[:, 1]*100, "o-", color="#697887", label="Entrenamiento")
    ax.plot(opciones_C, resultados_C[:, 2]*100, "s-", color="#007f89", label="Validación")
    ax.scatter(C_elegido, resultados_C[mejor, 2]*100, s=170, facecolors="none", edgecolors="#c96932", linewidths=2, label="Opción elegida")
    ax.scatter(C_vista, resultados_C[fila, 2]*100, s=70, color="#226b8a", label="C inspeccionado")
    ax.set(xscale="log", title="Elegir C sin consultar prueba", xlabel="C · menor valor = mayor regularización", ylabel="Exactitud equilibrada (%)", ylim=(50, 100))
    ax.legend()
    plt.show()
    print(f"Inspección: C={C_vista:g} · entrenamiento {resultados_C[fila,1]:.2%} · validación {resultados_C[fila,2]:.2%}")
    print(f"Modelo del informe: C={C_elegido:g}. Los botones sólo cambian esta vista.")
    return fig

C_control = widgets.ToggleButtons(options=opciones_C, value=C_elegido, description="Inspeccionar C",
                             style={"description_width": "initial"})
visor_candidatos = widgets.interactive(ver_candidatos, C_vista=C_control)
display(visor_candidatos)
\end{Codigo}

\Figura{paso-6e-1.png}{6E · Los candidatos comparten entrenamiento. La exactitud equilibrada de validación elige C; la prueba no participa en esa elección.}{94mm}

\begin{Salida}{6E · Resultado}
Inspección: C=0.1 · entrenamiento 95.37% · validación 95.05%
Modelo del informe: C=0.1. Los botones sólo cambian esta vista.
\end{Salida}

\Pagina{BLOQUE 7}{Ver una predicción}{}\label{bloque-7}

Aplicaremos el modelo a \textbf{una imagen de validación} y seguiremos la cuenta completa. \texttt{\small decision\_\allowbreak{}function} devuelve la suma z; \texttt{\small predict\_\allowbreak{}proba} devuelve puntuaciones por clase. Un \textbf{umbral} compara la puntuación de clase 1 con un número: \texttt{\small p >=\allowbreak{} umbral} produce la predicción 1; si no, produce 0. Referencia y predicción son dos datos distintos.


\needspace{10\baselineskip}\phantomsection\addcontentsline{toc}{subsection}{7A · Elegir el ejemplo de validación}\Sub{7A · Elegir el ejemplo de validación}\label{paso-7a}

\texttt{\small i\_\allowbreak{}val =\allowbreak{} 0} elige la primera imagen de validación conservada. \texttt{\small ids\_\allowbreak{}val[\allowbreak{}i\_\allowbreak{}val]\allowbreak{}} recuerda su posición en el archivo original, antes de quitar copias. El \texttt{\small assert} comprueba que el índice exista.

La misma imagen se seguirá en la suma, puntuación, predicción y comparación. Su referencia está en \texttt{\small y\_\allowbreak{}val[\allowbreak{}i\_\allowbreak{}val]\allowbreak{}}; no se pasa como argumento al pedir la predicción.


\begin{Codigo}{7A · Código}
i_val = 0
assert 0 <= i_val < len(X_val)
print("Validación limpia:", i_val, "· índice original:", ids_val[i_val])
\end{Codigo}

\begin{Salida}{7A · Resultado}
Validación limpia: 0 · índice original: 0
\end{Salida}

\needspace{10\baselineskip}\phantomsection\addcontentsline{toc}{subsection}{7B · Calcular puntuaciones de validación}\Sub{7B · Calcular puntuaciones de validación}\label{paso-7b}

\textbf{Aplicamos el modelo ya aprendido.} \texttt{\small decision\_\allowbreak{}function(\allowbreak{}X\_\allowbreak{}val)\allowbreak{}} calcula una suma z por imagen. Para la imagen elegida, \texttt{\small np.\allowbreak{}dot(\allowbreak{}X\_\allowbreak{}val[\allowbreak{}i\_\allowbreak{}val]\allowbreak{},\allowbreak{} modelo.\allowbreak{}coef\_\allowbreak{}[\allowbreak{}0]\allowbreak{})\allowbreak{} + modelo.\allowbreak{}intercept\_\allowbreak{}[\allowbreak{}0]\allowbreak{}} reproduce la cuenta con sus 784 intensidades, pesos y sesgo.

\texttt{\small predict\_\allowbreak{}proba(\allowbreak{}X\_\allowbreak{}val)\allowbreak{}} devuelve dos columnas de puntuaciones, una por clase. \texttt{\small classes\_\allowbreak{}} indica su orden; \texttt{\small [\allowbreak{}:\allowbreak{},\allowbreak{} columna\_\allowbreak{}1]\allowbreak{}} conserva la de clase 1. \texttt{\small allclose} comprueba que la cuenta manual coincide con z y que su sigmoide coincide con p. Lee los números impresos: \textbf{imagen concreta $\rightarrow$ z concreto $\rightarrow$ p concreto}. Ninguna de estas llamadas entrena otra vez.


\begin{Codigo}{7B · Código}
columna_1 = int(np.flatnonzero(modelo.classes_ == 1)[0])
z_val = modelo.decision_function(X_val)
p_val = modelo.predict_proba(X_val)[:, columna_1]
z_manual = float(np.dot(X_val[i_val], modelo.coef_[0]) + modelo.intercept_[0])
p_manual = float(1 / (1 + np.exp(-z_manual)))
assert np.allclose([z_manual, p_manual], [z_val[i_val], p_val[i_val]])
print("Orden de clases:", modelo.classes_, "· columna conservada:", columna_1)
print(f"Imagen de validación {i_val}: 784 intensidades × 784 pesos + sesgo → z = {z_manual:.6f}")
print(f"Sigmoide de z → {p_manual:.6f}; predict_proba para clase 1 → {p_val[i_val]:.6f}")
\end{Codigo}

\begin{Salida}{7B · Resultado}
Orden de clases: [0 1] · columna conservada: 1
Imagen de validación 0: 784 intensidades × 784 pesos + sesgo → z = 6.030946
Sigmoide de z → 0.997603; predict_proba para clase 1 → 0.997603
\end{Salida}

\needspace{10\baselineskip}\phantomsection\addcontentsline{toc}{subsection}{7C · Convertir una puntuación en decisión}\Sub{7C · Convertir una puntuación en decisión}\label{paso-7c}

\texttt{\small p\_\allowbreak{}ejemplo >=\allowbreak{} 0.\allowbreak{}5} pregunta si la puntuación alcanza el umbral. \texttt{\small int} convierte verdadero en 1 y falso en 0. \textbf{Puntuación} es un decimal; \textbf{predicción} es el rótulo decidido.

La salida muestra referencia y predicción por separado. Si coinciden, el ejemplo se cuenta como acierto. Una puntuación alta significa que el modelo favorece clase 1 bajo lo que aprendió; no demuestra una enfermedad ni una probabilidad clínica calibrada.


\begin{Codigo}{7C · Código}
p_ejemplo = float(p_val[i_val])
decision = int(p_ejemplo >= 0.5)
print(f"Puntuación {p_ejemplo:.6f} >= 0.5 → predicción {decision} = {nombres[decision]}")
print(f"Referencia publicada: {y_val[i_val]} = {nombres[y_val[i_val]]}")
print("¿Coinciden?", bool(decision == y_val[i_val]))
\end{Codigo}

\begin{Salida}{7C · Resultado}
Puntuación 0.997603 >= 0.5 → predicción 1 = Neumonía (etiqueta 1)
Referencia publicada: 1 = Neumonía (etiqueta 1)
¿Coinciden? True
\end{Salida}

\needspace{10\baselineskip}\phantomsection\addcontentsline{toc}{subsection}{7D · Relacionar entrada y salida}\Sub{7D · Relacionar entrada y salida}\label{paso-7d}

\textbf{Mueve Imagen} y observa la cuenta para otra radiografía. \textbf{Mueve Umbral} manteniendo la imagen: z y p permanecen iguales, pero la decisión puede cambiar.

\texttt{\small p\_\allowbreak{}val[\allowbreak{}indice\_\allowbreak{}validacion]\allowbreak{}} recupera la puntuación ya calculada. \texttt{\small int(\allowbreak{}puntuacion >=\allowbreak{} umbral\_\allowbreak{}vista)\allowbreak{}} decide 0/1. \texttt{\small imshow} muestra la entrada, \texttt{\small barh} dibuja p y \texttt{\small axvline} coloca el umbral. El recuadro une z, sigmoide y comparación; \texttt{\small set\_\allowbreak{}title} distingue referencia de predicción. Los controles son locales al visor: no alteran pesos ni el umbral del informe.


\begin{Codigo}{7D · Código}
def ver_prediccion(indice_validacion=0, umbral_vista=0.5):
    puntuacion = float(p_val[indice_validacion])
    prediccion = int(puntuacion >= umbral_vista)
    fig, ejes = plt.subplots(1, 3, figsize=(13, 4.4), layout="constrained", gridspec_kw={"width_ratios": [1, 1.2, 1.15]})
    ejes[0].imshow(imagenes_val[indice_validacion], cmap="gray", vmin=0, vmax=255, interpolation="nearest")
    ejes[0].set_title(f"Validación · índice original {ids_val[indice_validacion]}\nReferencia: {nombres[y_val[indice_validacion]]}")
    ejes[0].axis("off")
    ejes[1].barh(["Puntuación para etiqueta 1"], [puntuacion], height=0.3, color="#007f89")
    ejes[1].axvline(umbral_vista, color="#c96932", linestyle="--", label=f"Umbral {umbral_vista:.2f}")
    ejes[1].set(ylim=(-0.75, 0.75), xlim=(0, 1), xlabel="Salida estimada por el modelo", title=f"p={puntuacion:.3f} → {nombres[prediccion]}")
    ejes[1].legend(loc="lower right")
    ejes[2].axis("off")
    cuenta = (f"SUMA APRENDIDA\nz = {z_val[indice_validacion]:.3f}\n\n"
              f"SIGMOIDE\np = {puntuacion:.3f}\n\n"
              f"REGLA\np ≥ {umbral_vista:.2f} → {prediccion}\n"
              f"{nombres[prediccion]}")
    ejes[2].text(0.05, 0.95, cuenta, va="top", fontsize=12, transform=ejes[2].transAxes)
    ejes[2].set_title("De los píxeles a la decisión")
    plt.show()
    print(f"Referencia {y_val[indice_validacion]} · predicción {prediccion} · coinciden: {prediccion == y_val[indice_validacion]}")
    print("La puntuación no es una probabilidad clínica calibrada.")
    return fig

indice_control = widgets.IntSlider(value=i_val, min=0, max=len(X_val)-1, step=1,
                                   description="Imagen", continuous_update=False)
decision_control = widgets.FloatSlider(value=0.5, min=0.0, max=1.0, step=0.01,
                                     description="Umbral", continuous_update=False)
visor_prediccion = widgets.interactive(ver_prediccion, indice_validacion=indice_control,
                                     umbral_vista=decision_control)
display(visor_prediccion)
\end{Codigo}

\Figura{paso-7d-1.png}{7D · Una imagen pasa por la suma con pesos, la sigmoide y la comparación con el umbral. La referencia permite comprobar la predicción.}{94mm}

\begin{Salida}{7D · Resultado}
Referencia 1 · predicción 1 · coinciden: True
La puntuación no es una probabilidad clínica calibrada.
\end{Salida}

\Pagina{BLOQUE 8}{Mover el umbral}{}\label{bloque-8}

Las puntuaciones ya están calculadas. Ahora cambiaremos sólo la regla que las convierte en etiquetas. Con p = 0.40, el umbral 0.50 produce 0 y el umbral 0.30 produce 1. Al bajar el umbral habrá más predicciones 1; algunas corregirán errores y otras crearán errores nuevos.


\Sub{Leer una matriz con cuatro casos ficticios}\Figura{cuatro_casos_confusion.png}{Cuatro casos ficticios con puntuaciones inventadas. La matriz cuenta una vez cada caso: no describe pacientes ni resultados del conjunto médico.}{76mm}
\begin{enumerate}
\item \textbf{Decide con el umbral.} Las puntuaciones \([0,20;0,80;0,30;0,70]\) producen predicciones \([0;1;0;1]\) al aplicar \(s\geq0,5\).
\item \textbf{Busca el cruce.} La referencia elige la fila; la predicción elige la columna. C tiene referencia 1 y predicción 0: cae en falso negativo, abajo a la izquierda.
\item \textbf{Cuenta y comprueba.} Cada casilla contiene un caso. La suma es cuatro: dos coincidencias en la diagonal y dos discrepancias fuera de ella.
\end{enumerate}
\textbf{Positivo/negativo} nombra la predicción 1/0. \textbf{Verdadero/falso} indica si coincide con la referencia. La referencia se consulta para evaluar, no se entrega como entrada al predecir. Estos términos no indican gravedad.
\begin{Codigo}{La misma cuenta con Python · ejemplo ficticio}
y_ejemplo = np.array([0, 0, 1, 1])
s_ejemplo = np.array([0.20, 0.80, 0.30, 0.70])
pred_ejemplo = (s_ejemplo >= 0.5).astype(int)
confusion_matrix(y_ejemplo, pred_ejemplo, labels=[0, 1])
# Salida por filas: [[1, 1], [1, 1]]
\end{Codigo}
\begin{Practica}{Ahora baja el umbral a 0,3}
C pasa de FN a VP: 0,30 cumple \texttt{>= 0.3}, que incluye la igualdad. La matriz queda \texttt{[[1, 1], [0, 2]]}. Los otros casos no cambian y la suma sigue siendo cuatro. No han cambiado las puntuaciones ni los pesos.
\end{Practica}



\needspace{10\baselineskip}\phantomsection\addcontentsline{toc}{subsection}{8A · Fijar los umbrales de demostración}\Sub{8A · Fijar los umbrales de demostración}\label{paso-8a}

Comparamos las reglas 0.3, 0.5 y 0.7 usando las mismas puntuaciones de validación. \texttt{\small umbrales\_\allowbreak{}demo} es una lista, no tres modelos nuevos.

Necesitamos nombrar cuatro casos: \textbf{VN}, referencia 0 y predicción 0; \textbf{FP}, referencia 0 y predicción 1; \textbf{FN}, referencia 1 y predicción 0; \textbf{VP}, referencia 1 y predicción 1. «Verdadero» indica coincidencia; «falso», discrepancia.


\begin{Codigo}{8A · Código}
umbrales_demo = [0.3, 0.5, 0.7]
print("Umbrales ilustrativos de validación:", umbrales_demo)
\end{Codigo}

\begin{Salida}{8A · Resultado}
Umbrales ilustrativos de validación: [0.3, 0.5, 0.7]
\end{Salida}

\needspace{10\baselineskip}\phantomsection\addcontentsline{toc}{subsection}{8B · Contar las tres matrices}\Sub{8B · Contar las tres matrices}\label{paso-8b}

\texttt{\small p\_\allowbreak{}val >=\allowbreak{} umbral} compara todas las puntuaciones y \texttt{\small .\allowbreak{}astype(\allowbreak{}int)\allowbreak{}} convierte las respuestas en etiquetas. \texttt{\small confusion\_\allowbreak{}matrix(\allowbreak{}y\_\allowbreak{}val,\allowbreak{} pred,\allowbreak{} labels=\allowbreak{}[\allowbreak{}0,\allowbreak{} 1]\allowbreak{})\allowbreak{}} cuenta las combinaciones.

Lee cada matriz por filas: \texttt{\small [\allowbreak{}[\allowbreak{}VN,\allowbreak{} FP]\allowbreak{},\allowbreak{} [\allowbreak{}FN,\allowbreak{} VP]\allowbreak{}]\allowbreak{}}. La primera fila contiene referencias 0; la segunda, referencias 1. Las columnas contienen predicciones 0 y 1. El mensaje escribe los nombres y totales para no depender de memorizar posiciones. \texttt{\small append} conserva cada matriz de la comparación.


\begin{Codigo}{8B · Código}
matrices_val = []
for umbral in umbrales_demo:
    pred = (p_val >= umbral).astype(int)
    matriz = confusion_matrix(y_val, pred, labels=[0, 1])
    matrices_val.append(matriz)
    print(f"Umbral {umbral}: VN={matriz[0,0]}, FP={matriz[0,1]}, FN={matriz[1,0]}, VP={matriz[1,1]} · total {matriz.sum()}")
\end{Codigo}

\begin{Salida}{8B · Resultado}
Umbral 0.3: VN=122, FP=13, FN=9, VP=372 · total 516
Umbral 0.5: VN=128, FP=7, FN=18, VP=363 · total 516
Umbral 0.7: VN=131, FP=4, FN=53, VP=328 · total 516
\end{Salida}

\needspace{10\baselineskip}\phantomsection\addcontentsline{toc}{subsection}{8C · Comparar visualmente los errores}\Sub{8C · Comparar visualmente los errores}\label{paso-8c}

\textbf{Mueve Umbral} y lee los cuatro recuadros. Cada uno combina una referencia, una predicción y un número de imágenes. En 0.5, la salida guardada cuenta 128 VN, 7 FP, 18 FN y 363 VP; suman 516.

La función recalcula etiquetas con \texttt{\small >=\allowbreak{}} y conteos con \texttt{\small confusion\_\allowbreak{}matrix}. \texttt{\small imshow} colorea los conteos; \texttt{\small np.\allowbreak{}ndenumerate} recorre fila, columna y valor para escribir sigla y cantidad. Cambiar el umbral redistribuye casos entre casillas; no cambia las imágenes, referencias ni puntuaciones. El total se conserva y el informe mantiene su decisión fija.


\begin{Codigo}{8C · Código}
def ver_umbral(umbral_validacion=0.5):
    prediccion = (p_val >= umbral_validacion).astype(int)
    matriz = confusion_matrix(y_val, prediccion, labels=[0, 1])
    fig, ax = plt.subplots(figsize=(7.8, 5.6), layout="constrained")
    ax.imshow(matriz, cmap="Blues", vmin=0, vmax=len(y_val))
    siglas = np.array([["VN", "FP"], ["FN", "VP"]])
    for (fila, columna), valor in np.ndenumerate(matriz):
        ax.text(columna, fila, f"{siglas[fila,columna]}\n{valor}", ha="center", va="center", fontsize=22,
                color="white" if valor > len(y_val)/2 else "#102c40")
    ax.set(xticks=[0, 1], xticklabels=["Normal · 0", "Neumonía · 1"],
           yticks=[0, 1], yticklabels=["Normal · 0", "Neumonía · 1"],
           xlabel="Predicción del modelo", ylabel="Etiqueta de referencia",
           title=f"Validación · umbral {umbral_validacion:.2f}\nFP={matriz[0,1]} · FN={matriz[1,0]}")
    plt.show()
    print(f"VN={matriz[0,0]}, FP={matriz[0,1]}, FN={matriz[1,0]}, VP={matriz[1,1]}")
    return fig

umbral_control = widgets.FloatSlider(value=0.5, min=0.0, max=1.0, step=0.01,
                                    description="Umbral", continuous_update=False)
visor_umbral = widgets.interactive(ver_umbral, umbral_validacion=umbral_control)
display(visor_umbral)
\end{Codigo}

\Figura{paso-8c-1.png}{8C · La matriz de validación cuenta cada combinación de referencia y predicción. Mover el umbral cambia las decisiones, no las puntuaciones.}{94mm}

\begin{Salida}{8C · Resultado}
VN=128, FP=7, FN=18, VP=363
\end{Salida}

\needspace{10\baselineskip}\phantomsection\addcontentsline{toc}{subsection}{8D · Confirmar las decisiones del informe}\Sub{8D · Confirmar las decisiones del informe}\label{paso-8d}

Escribimos \texttt{\small UMBRAL\_\allowbreak{}FINAL =\allowbreak{} 0.\allowbreak{}5}: ésta es la regla que se aplicará a prueba. No la elegimos mirando el resultado de prueba. Los controles de exploración usan variables locales y no cambian \texttt{\small UMBRAL\_\allowbreak{}FINAL}.

Dejamos fijados \textbf{modelo elegido por validación + umbral 0.5}. A partir de aquí, la finalidad es informar su rendimiento, no buscar una cifra de prueba más favorable.


\begin{Codigo}{8D · Código}
UMBRAL_FINAL = 0.5
print(f"Decisiones fijadas: C={C_elegido:g}; umbral={UMBRAL_FINAL}; no se reajusta tras abrir prueba.")
\end{Codigo}

\begin{Salida}{8D · Resultado}
Decisiones fijadas: C=0.1; umbral=0.5; no se reajusta tras abrir prueba.
\end{Salida}

\Pagina{BLOQUE 9}{Evaluar con decisiones fijas}{}\label{bloque-9}

Usamos las 624 imágenes de prueba por primera vez para evaluar el modelo elegido. Una \textbf{matriz de confusión} cuenta cuatro combinaciones: referencia 0 o 1 frente a predicción 0 o 1. La diagonal contiene coincidencias; las otras casillas contienen discrepancias.


\needspace{10\baselineskip}\phantomsection\addcontentsline{toc}{subsection}{9A · Abrir la partición de prueba}\Sub{9A · Abrir la partición de prueba}\label{paso-9a}

Abrimos ahora las imágenes y referencias reservadas para prueba. Las claves \texttt{\small test\_\allowbreak{}images} y \texttt{\small test\_\allowbreak{}labels} entregan sus matrices y \texttt{\small .\allowbreak{}ravel(\allowbreak{})\allowbreak{}} deja una etiqueta por imagen.

Los \texttt{\small assert} comprueban 624 imágenes de 28 $\times$ 28 y \texttt{\small bincount} verifica 234 referencias 0 y 390 referencias 1. Usaremos esas etiquetas para evaluar después de predecir; no se pasarán al modelo para ajustar pesos.


\begin{Codigo}{9A · Código}
imagenes_test = paquete["test_images"]
y_test = paquete["test_labels"].ravel()
assert imagenes_test.shape == (624, 28, 28)
assert np.array_equal(np.bincount(y_test), [234, 390])
print("Prueba oficial:", imagenes_test.shape, "· etiquetas:", np.bincount(y_test))
\end{Codigo}

\begin{Salida}{9A · Resultado}
Prueba oficial: (624, 28, 28) · etiquetas: [234 390]
\end{Salida}

\needspace{10\baselineskip}\phantomsection\addcontentsline{toc}{subsection}{9B · Comprobar coincidencias exactas}\Sub{9B · Comprobar coincidencias exactas}\label{paso-9b}

Verificamos que ninguna imagen de prueba sea una copia exacta de las usadas para aprender o validar. \texttt{\small |} une los conjuntos de bytes de entrenamiento y validación; \texttt{\small not in} comprueba cada imagen de prueba y \texttt{\small all} exige que todas pasen.

No se elimina ninguna imagen de prueba. La ausencia de copias idénticas no demuestra que sean pacientes diferentes; esa limitación seguirá en el informe.


\begin{Codigo}{9B · Código}
huellas_val = {imagen.tobytes() for imagen in imagenes_val}
assert all(imagen.tobytes() not in huellas_train | huellas_val for imagen in imagenes_test)
print("No hay coincidencias exactas con train/val; esto no certifica independencia por paciente.")
\end{Codigo}

\begin{Salida}{9B · Resultado}
No hay coincidencias exactas con train/val; esto no certifica independencia por paciente.
\end{Salida}

\needspace{10\baselineskip}\phantomsection\addcontentsline{toc}{subsection}{9C · Preparar las características de prueba}\Sub{9C · Preparar las características de prueba}\label{paso-9c}

Preparamos la entrada con la misma regla para todas las particiones: \texttt{\small reshape} deja una imagen por fila, \texttt{\small astype} permite decimales y \texttt{\small /\allowbreak{}255.\allowbreak{}0} conserva la escala.

La salida debe ser \texttt{\small (\allowbreak{}624,\allowbreak{} 784)\allowbreak{}}: 624 imágenes que el modelo recibirá, cada una con 784 intensidades. Cambiar la regla de preparación aquí haría injusta la comparación con lo aprendido.


\begin{Codigo}{9C · Código}
X_test = imagenes_test.reshape(len(imagenes_test), -1).astype(np.float64) / 255.0
print("X_test:", X_test.shape)
\end{Codigo}

\begin{Salida}{9C · Resultado}
X_test: (624, 784)
\end{Salida}

\needspace{10\baselineskip}\phantomsection\addcontentsline{toc}{subsection}{9D · Predecir con las decisiones fijadas}\Sub{9D · Predecir con las decisiones fijadas}\label{paso-9d}

\texttt{\small modelo.\allowbreak{}predict\_\allowbreak{}proba(\allowbreak{}X\_\allowbreak{}test)\allowbreak{}} aplica los pesos fijados y devuelve puntuaciones. Elegimos la columna de clase 1; la comparación \texttt{\small >=\allowbreak{} UMBRAL\_\allowbreak{}FINAL} produce 624 etiquetas.

No hay llamada a \texttt{\small fit} ni se entrega \texttt{\small y\_\allowbreak{}test} al modelo. Las decisiones quedan guardadas en \texttt{\small pred\_\allowbreak{}test}. Las referencias se usarán para contar coincidencias, no para cambiar esas decisiones.


\begin{Codigo}{9D · Código}
p_test = modelo.predict_proba(X_test)[:, columna_1]
pred_test = (p_test >= UMBRAL_FINAL).astype(int)
print("Predicciones de prueba:", len(pred_test), "· umbral fijo:", UMBRAL_FINAL)
\end{Codigo}

\begin{Salida}{9D · Resultado}
Predicciones de prueba: 624 · umbral fijo: 0.5
\end{Salida}

\needspace{10\baselineskip}\phantomsection\addcontentsline{toc}{subsection}{9E · Construir la matriz de prueba}\Sub{9E · Construir la matriz de prueba}\label{paso-9e}

\texttt{\small confusion\_\allowbreak{}matrix(\allowbreak{}y\_\allowbreak{}test,\allowbreak{} pred\_\allowbreak{}test,\allowbreak{} labels=\allowbreak{}[\allowbreak{}0,\allowbreak{} 1]\allowbreak{})\allowbreak{}} compara referencia con predicción. \texttt{\small .\allowbreak{}ravel(\allowbreak{})\allowbreak{}} lee por filas y asignamos sus valores a \texttt{\small vn,\allowbreak{} fp,\allowbreak{} fn,\allowbreak{} vp}.

En la salida guardada: \textbf{158 + 76 = 234} referencias 0 y \textbf{14 + 376 = 390} referencias 1. Hay 158 + 376 coincidencias y 76 + 14 discrepancias. Los mensajes del código realizan estas cuentas con tus resultados.


\begin{Codigo}{9E · Código}
matriz_test = confusion_matrix(y_test, pred_test, labels=[0, 1])
vn, fp, fn, vp = [int(v) for v in matriz_test.ravel()]
print("Matriz:", matriz_test.tolist(), "· VN, FP, FN, VP:", vn, fp, fn, vp)
print(f"Referencias 0: {vn}+{fp}={vn+fp}; referencias 1: {fn}+{vp}={fn+vp}.")
print(f"Coincidencias: {vn+vp}; discrepancias: {fp+fn}; total: {vn+fp+fn+vp}.")
\end{Codigo}

\begin{Salida}{9E · Resultado}
Matriz: [[158, 76], [14, 376]] · VN, FP, FN, VP: 158 76 14 376
Referencias 0: 158+76=234; referencias 1: 14+376=390.
Coincidencias: 534; discrepancias: 90; total: 624.
\end{Salida}

\needspace{10\baselineskip}\phantomsection\addcontentsline{toc}{subsection}{9F · Ver dónde se cuentan los resultados}\Sub{9F · Ver dónde se cuentan los resultados}\label{paso-9f}

Para leer cualquier casilla, empieza por la \textbf{fila de referencia} y después por la \textbf{columna de predicción}. Por ejemplo, fila «Normal · 0» y columna «Neumonía · 1» cuenta falsos positivos. La diagonal VN/VP contiene aciertos respecto a las etiquetas publicadas.

\texttt{\small imshow} convierte conteos en color; \texttt{\small ndenumerate} proporciona posición y valor para \texttt{\small text}. Las etiquetas de ejes explican qué representa cada dirección. El título conserva C y umbral, porque son parte de la evaluación que estás viendo.


\begin{Codigo}{9F · Código}
fig, ax = plt.subplots(figsize=(7.2, 5), layout="constrained")
ax.imshow(matriz_test, cmap="Blues", vmin=0, vmax=max(matriz_test.max(), 1))
siglas = np.array([["VN", "FP"], ["FN", "VP"]])
for (fila, columna), valor in np.ndenumerate(matriz_test):
    ax.text(columna, fila, f"{siglas[fila,columna]}\n{valor}", ha="center", va="center", fontsize=22, color="white" if valor > matriz_test.max()/2 else "#102c40")
ax.set(xticks=[0, 1], xticklabels=["Normal · 0", "Neumonía · 1"], yticks=[0, 1], yticklabels=["Normal · 0", "Neumonía · 1"], xlabel="Predicción del modelo", ylabel="Etiqueta de referencia", title=f"Prueba oficial · 624 imágenes\nC={C_elegido:g} · umbral {UMBRAL_FINAL}")
plt.show()
print(f"VN={vn}, FP={fp}, FN={fn}, VP={vp} · total={vn+fp+fn+vp}")
\end{Codigo}

\Figura{paso-9f-1.png}{9F · Las 624 imágenes de prueba se reparten entre VN=158, FP=76, FN=14 y VP=376. Filas: referencia; columnas: predicción.}{94mm}

\begin{Salida}{9F · Resultado}
VN=158, FP=76, FN=14, VP=376 · total=624
\end{Salida}

\Pagina{BLOQUE 10}{Entender cada medida}{}\label{bloque-10}

Cada medida responde una pregunta usando los mismos cuatro conteos. En la salida guardada hay 158 verdaderos negativos (VN), 76 falsos positivos (FP), 14 falsos negativos (FN) y 376 verdaderos positivos (VP). \textbf{Positivo significa etiqueta 1}, no una valoración de si el resultado es bueno. Las cuentas del código utilizan tus propios resultados.


\needspace{10\baselineskip}\phantomsection\addcontentsline{toc}{subsection}{10A · Calcular sensibilidad}\Sub{10A · Calcular sensibilidad}\label{paso-10a}

\textbf{Sensibilidad:} de todas las referencias 1, ¿qué parte predijimos como 1? El denominador \texttt{\small vp + fn} incluye todas las referencias 1; el numerador \texttt{\small vp} sólo las reconocidas.

La cuenta guardada es \textbf{376 / (376 + 14) = 96.41\%}. No es el porcentaje de aciertos sobre todas las imágenes: sólo responde por la clase 1. La división del código y su mensaje muestran tu numerador y denominador.


\begin{Codigo}{10A · Código}
sensibilidad = vp / (vp + fn)
print(f"Referencias 1 reconocidas: VP/(VP+FN) = {vp}/({vp}+{fn}) = {sensibilidad:.2%}")
\end{Codigo}

\begin{Salida}{10A · Resultado}
Referencias 1 reconocidas: VP/(VP+FN) = 376/(376+14) = 96.41%
\end{Salida}

\needspace{10\baselineskip}\phantomsection\addcontentsline{toc}{subsection}{10B · Calcular especificidad}\Sub{10B · Calcular especificidad}\label{paso-10b}

\textbf{Especificidad:} de todas las referencias 0, ¿qué parte predijimos como 0? \texttt{\small vn + fp} cuenta todas las referencias 0; \texttt{\small vn} cuenta las reconocidas.

La cuenta guardada es \textbf{158 / (158 + 76) = 67.52\%}. Compárala con la sensibilidad: el modelo reconoce mejor una clase que la otra en este conjunto. La escala de ambas va de 0 a 100\%.


\begin{Codigo}{10B · Código}
especificidad = vn / (vn + fp)
print(f"Referencias 0 reconocidas: VN/(VN+FP) = {vn}/({vn}+{fp}) = {especificidad:.2%}")
\end{Codigo}

\begin{Salida}{10B · Resultado}
Referencias 0 reconocidas: VN/(VN+FP) = 158/(158+76) = 67.52%
\end{Salida}

\needspace{10\baselineskip}\phantomsection\addcontentsline{toc}{subsection}{10C · Calcular exactitud}\Sub{10C · Calcular exactitud}\label{paso-10c}

\textbf{Exactitud:} de todas las imágenes, ¿cuántas predicciones coincidieron con su referencia? \texttt{\small vp + vn} suma la diagonal y \texttt{\small len(\allowbreak{}y\_\allowbreak{}test)\allowbreak{}} cuenta todas las imágenes.

La salida guardada es \textbf{(376 + 158) / 624 = 85.58\%}. Cada imagen pesa igual; por eso la clase con más imágenes influye más en esta cifra.


\begin{Codigo}{10C · Código}
exactitud = (vp + vn) / len(y_test)
print(f"Coincidencias entre todas: (VP+VN)/N = ({vp}+{vn})/{len(y_test)} = {exactitud:.2%}")
\end{Codigo}

\begin{Salida}{10C · Resultado}
Coincidencias entre todas: (VP+VN)/N = (376+158)/624 = 85.58%
\end{Salida}

\needspace{10\baselineskip}\phantomsection\addcontentsline{toc}{subsection}{10D · Calcular exactitud equilibrada}\Sub{10D · Calcular exactitud equilibrada}\label{paso-10d}

\textbf{Exactitud equilibrada:} promediamos lo reconocido en cada clase: \texttt{\small (\allowbreak{}sensibilidad + especificidad)\allowbreak{} /\allowbreak{} 2}.

Con las cifras guardadas, \texttt{\small (\allowbreak{}96.\allowbreak{}41\% + 67.\allowbreak{}52\%)\allowbreak{}/\allowbreak{}2} es aproximadamente \textbf{81.97\%}. Cada clase aporta la mitad del resultado, aunque sus tamaños sean diferentes. Es una pregunta distinta de la exactitud; no se deben comparar sus números como si fueran la misma medida.


\begin{Codigo}{10D · Código}
equilibrada = (sensibilidad + especificidad) / 2
print(f"Promedio por clase: ({sensibilidad:.2%} + {especificidad:.2%})/2 = {equilibrada:.2%}")
\end{Codigo}

\begin{Salida}{10D · Resultado}
Promedio por clase: (96.41% + 67.52%)/2 = 81.97%
\end{Salida}

\needspace{10\baselineskip}\phantomsection\addcontentsline{toc}{subsection}{10E · Aplicar también la referencia a prueba}\Sub{10E · Aplicar también la referencia a prueba}\label{paso-10e}

Aplicamos la regla mayoritaria a las mismas imágenes de prueba. \texttt{\small .\allowbreak{}predict} no cambia esa regla. \texttt{\small confusion\_\allowbreak{}matrix} usa otra vez referencia en filas y predicción en columnas.

La regla predice siempre clase 1: reconoce las 390 referencias 1 y falla las 234 referencias 0. Puede acertar 62.5\% del total sin reconocer ninguna imagen de clase 0. Guardamos sus conteos para comparar con el modelo aprendido.


\begin{Codigo}{10E · Código}
pred_base_test = referencia.predict(X_test)
bn, bf, bfn, bp = confusion_matrix(y_test, pred_base_test, labels=[0, 1]).ravel()
print("Referencia mayoritaria · VN, FP, FN, VP:", bn, bf, bfn, bp)
\end{Codigo}

\begin{Salida}{10E · Resultado}
Referencia mayoritaria · VN, FP, FN, VP: 0 234 0 390
\end{Salida}

\needspace{10\baselineskip}\phantomsection\addcontentsline{toc}{subsection}{10F · Preparar las medidas para comparar}\Sub{10F · Preparar las medidas para comparar}\label{paso-10f}

Reunimos las cuatro medidas en un orden compartido por ambos modelos. Las listas \texttt{\small metricas} y \texttt{\small base\_\allowbreak{}test} deben conservar el mismo orden que \texttt{\small etiquetas\_\allowbreak{}metricas}; así cada par de barras compara la misma pregunta.

\texttt{\small np.\allowbreak{}arange(\allowbreak{}4)\allowbreak{}} crea sus posiciones. Los saltos de línea dividen rótulos largos para que se lean completos. La salida presenta nombres y porcentajes, no una lista de decimales sin identificar.


\begin{Codigo}{10F · Código}
metricas = [exactitud, equilibrada, sensibilidad, especificidad]
base_test = [accuracy_score(y_test, pred_base_test), balanced_accuracy_score(y_test, pred_base_test), bp/(bp+bfn), bn/(bn+bf)]
etiquetas_metricas = ["Exactitud", "Exactitud\nequilibrada", "Sensibilidad\npara etiqueta 1", "Especificidad\npara etiqueta 0"]
x = np.arange(4)
for nombre, valor, base in zip(etiquetas_metricas, metricas, base_test):
    print(f"{nombre.replace(chr(10), ' ')}: modelo {valor:.2%} · regla mayoritaria {base:.2%}")
\end{Codigo}

\begin{Salida}{10F · Resultado}
Exactitud: modelo 85.58% · regla mayoritaria 62.50%
Exactitud equilibrada: modelo 81.97% · regla mayoritaria 50.00%
Sensibilidad para etiqueta 1: modelo 96.41% · regla mayoritaria 100.00%
Especificidad para etiqueta 0: modelo 67.52% · regla mayoritaria 0.00%
\end{Salida}

\needspace{10\baselineskip}\phantomsection\addcontentsline{toc}{subsection}{10G · Comparar y explicar los denominadores}\Sub{10G · Comparar y explicar los denominadores}\label{paso-10g}

Compara \textbf{las dos barras de una misma medida}: gris es la regla constante y turquesa la regresión logística. Una sensibilidad alta con especificidad baja significa que se reconocen muchas referencias 1 pero se confunden muchas referencias 0.

\texttt{\small x-0.\allowbreak{}19} y \texttt{\small x+0.\allowbreak{}19} colocan las barras lado a lado; \texttt{\small bar\_\allowbreak{}label} escribe porcentajes y \texttt{\small legend} identifica los modelos. \texttt{\small print} muestra los numeradores y denominadores. Superar la referencia en este conjunto no establece utilidad clínica ni rendimiento en otro hospital.


\begin{Codigo}{10G · Código}
fig, ax = plt.subplots(figsize=(11, 4.6), layout="constrained")
b1 = ax.bar(x-0.19, np.array(base_test)*100, 0.38, label="Regla mayoritaria", color="#9ba8b0")
b2 = ax.bar(x+0.19, np.array(metricas)*100, 0.38, label="Regresión logística", color="#007f89")
ax.bar_label(b1, labels=[f"{v:.1%}" for v in base_test], padding=3, fontsize=9)
ax.bar_label(b2, labels=[f"{v:.1%}" for v in metricas], padding=3, fontsize=9)
ax.set(xticks=x, xticklabels=etiquetas_metricas, ylim=(0, 125), ylabel="Porcentaje", title="Prueba · cuatro preguntas distintas")
ax.legend(loc="upper center", ncols=2)
plt.show()
print(f"Sensibilidad: VP/(VP+FN) = {vp}/{vp+fn} = {sensibilidad:.2%}")
print(f"Especificidad: VN/(VN+FP) = {vn}/{vn+fp} = {especificidad:.2%}")
print(f"Exactitud: (VP+VN)/N = {vp+vn}/{len(y_test)} = {exactitud:.2%}")
print(f"Equilibrada: (sensibilidad+especificidad)/2 = {equilibrada:.2%}")
\end{Codigo}

\Figura{paso-10g-1.png}{10G · Cada medida utiliza un denominador diferente. La exactitud equilibrada promedia sensibilidad y especificidad, dando el mismo peso a ambas clases.}{94mm}

\begin{Salida}{10G · Resultado}
Sensibilidad: VP/(VP+FN) = 376/390 = 96.41%
Especificidad: VN/(VN+FP) = 158/234 = 67.52%
Exactitud: (VP+VN)/N = 534/624 = 85.58%
Equilibrada: (sensibilidad+especificidad)/2 = 81.97%
\end{Salida}

\Pagina{BLOQUE 11}{Mirar aciertos y errores}{}\label{bloque-11}

Vamos a ver imágenes que pertenecen a cada casilla. Un \textbf{falso positivo} tiene referencia 0 y predicción 1; un \textbf{falso negativo}, referencia 1 y predicción 0. Las imágenes ilustran esas combinaciones. Mirarlas no demuestra qué rasgo usó el modelo ni permite reinterpretar médicamente la etiqueta.


\needspace{10\baselineskip}\phantomsection\addcontentsline{toc}{subsection}{11A · Definir las cuatro casillas}\Sub{11A · Definir las cuatro casillas}\label{paso-11a}

Las cuatro tuplas enumeran las casillas: \texttt{\small (\allowbreak{}0,\allowbreak{}0)\allowbreak{}} es VN, \texttt{\small (\allowbreak{}0,\allowbreak{}1)\allowbreak{}} es FP, \texttt{\small (\allowbreak{}1,\allowbreak{}0)\allowbreak{}} es FN y \texttt{\small (\allowbreak{}1,\allowbreak{}1)\allowbreak{}} es VP. El primer número es siempre referencia; el segundo, predicción.

Cada tupla añade un rótulo legible. \texttt{\small titulo[\allowbreak{}:\allowbreak{}2]\allowbreak{}} toma la sigla para usarla como clave del diccionario de ejemplos.


\begin{Codigo}{11A · Código}
casos = [(0, 0, "VN · coincide en 0"), (0, 1, "FP · predice 1, referencia 0"), (1, 0, "FN · predice 0, referencia 1"), (1, 1, "VP · coincide en 1")]
print("Casillas:", [titulo[:2] for _, _, titulo in casos])
\end{Codigo}

\begin{Salida}{11A · Resultado}
Casillas: ['VN', 'FP', 'FN', 'VP']
\end{Salida}

\needspace{10\baselineskip}\phantomsection\addcontentsline{toc}{subsection}{11B · Elegir el primer ejemplo de cada casilla}\Sub{11B · Elegir el primer ejemplo de cada casilla}\label{paso-11b}

Buscamos el primer ejemplo de cada combinación. \texttt{\small \&} exige que se cumplan a la vez la referencia y la predicción; \texttt{\small flatnonzero} obtiene las posiciones que cumplen.

Si la casilla no está vacía, \texttt{\small candidatos[\allowbreak{}0]\allowbreak{}} toma su primera posición. Esta regla evita escoger a mano una imagen que parezca especialmente convincente. Si una casilla no tiene ejemplos, el visor lo indicará.


\begin{Codigo}{11B · Código}
indices_ejemplos = {}
for real, predicho, titulo in casos:
    candidatos = np.flatnonzero((y_test == real) & (pred_test == predicho))
    if len(candidatos):
        indices_ejemplos[titulo[:2]] = int(candidatos[0])
print("Primeros índices disponibles:", indices_ejemplos)
\end{Codigo}

\begin{Salida}{11B · Resultado}
Primeros índices disponibles: {'VN': 3, 'FP': 1, 'FN': 110, 'VP': 0}
\end{Salida}

\needspace{10\baselineskip}\phantomsection\addcontentsline{toc}{subsection}{11C · Mostrar los ejemplos seleccionados}\Sub{11C · Mostrar los ejemplos seleccionados}\label{paso-11c}

\textbf{Pulsa Todas, VN, FP, FN o VP} para comparar o ampliar los ejemplos fijados. Cada título muestra referencia y predicción por separado, junto con su puntuación.

\texttt{\small get} recupera el índice guardado. \texttt{\small imshow} dibuja sus píxeles y \texttt{\small set\_\allowbreak{}title} describe por qué cae en esa casilla. \texttt{\small ToggleButtons} ofrece botones de selección. Cambiar la vista no sustituye ejemplos, no vuelve a evaluar prueba y no modifica el modelo. Un error respecto a la etiqueta no permite explicar por sí solo qué región de la imagen lo causó.


\begin{Codigo}{11C · Código}
def ver_casilla(casilla="Todas"):
    seleccion = casos if casilla == "Todas" else [c for c in casos if c[2][:2] == casilla]
    filas, columnas = (2, 2) if casilla == "Todas" else (1, 1)
    fig, ejes = plt.subplots(filas, columnas, figsize=(9, 7.4) if casilla == "Todas" else (5, 5), layout="constrained")
    for ax, (real, predicho, titulo) in zip(np.atleast_1d(ejes).ravel(), seleccion):
        i = indices_ejemplos.get(titulo[:2])
        if i is not None:
            ax.imshow(imagenes_test[i], cmap="gray", vmin=0, vmax=255, interpolation="nearest")
            ax.set_title(f"{titulo}\nTest[{i}] · puntuación 1: {p_test[i]:.3f}", color="#007f89" if real == predicho else "#a24c23")
        else:
            ax.text(0.5, 0.5, "Sin ejemplos en esta casilla", ha="center", transform=ax.transAxes)
            ax.set_title(titulo)
        ax.axis("off")
    plt.show()
    print("Criterio: primer índice de prueba en cada casilla; no selección por apariencia ni puntuación.")
    return fig

casilla_control = widgets.ToggleButtons(options=["Todas", "VN", "FP", "FN", "VP"],
                                   value="Todas", description="Casilla")
visor_casilla = widgets.interactive(ver_casilla, casilla=casilla_control)
display(visor_casilla)
\end{Codigo}

\Figura{paso-11c-1.png}{11C · Primer ejemplo disponible de cada casilla. La selección usa los índices, no el aspecto de la radiografía.}{94mm}

\begin{Salida}{11C · Resultado}
Criterio: primer índice de prueba en cada casilla; no selección por apariencia ni puntuación.
\end{Salida}

\Pagina{BLOQUE 12}{Guardar el experimento}{}\label{bloque-12}

Para documentar el resultado guardaremos una ficha: datos usados, decisiones, conteos, medidas y versiones. \textbf{JSON} es un archivo de texto organizado por nombres y valores; no contiene el modelo entrenado ni sustituye el notebook. Permite revisar cómo se obtuvo la evaluación.


\needspace{10\baselineskip}\phantomsection\addcontentsline{toc}{subsection}{12A · Registrar los datos del experimento}\Sub{12A · Registrar los datos del experimento}\label{paso-12a}

Creamos \texttt{\small informe}, un \textbf{diccionario}: una colección de nombres asociados a valores. Por ejemplo, una clave de tamaño conserva cuántas imágenes se usaron.

Guardamos procedencia, huella del archivo, tamaños originales y efectivos, y exclusiones. Las llaves interiores agrupan datos de una misma categoría. Esta ficha explica qué datos sostienen las cifras.


\begin{Codigo}{12A · Código}
informe = {
    "dataset": "PneumoniaMNIST 28x28 · Zenodo 10519652",
    "sha256": sha_esperada,
    "particiones_originales": {"train": 4708, "val": 524, "test": 624},
    "particiones_utilizadas": {"train": len(y_train), "val": len(y_val), "test": len(y_test)},
    "exclusiones": "18 copias exactas adicionales en train; 8 val idénticas a train; test oficial intacto",
}
print("Ficha iniciada:", informe["particiones_utilizadas"])
\end{Codigo}

\begin{Salida}{12A · Resultado}
Ficha iniciada: {'train': 4690, 'val': 516, 'test': 624}
\end{Salida}

\needspace{10\baselineskip}\phantomsection\addcontentsline{toc}{subsection}{12B · Registrar modelo y decisiones}\Sub{12B · Registrar modelo y decisiones}\label{paso-12b}

\texttt{\small .\allowbreak{}update(\allowbreak{}\{.\allowbreak{}.\allowbreak{}.\allowbreak{}\})\allowbreak{}} añade decisiones al diccionario existente: tipo de modelo, C elegido, opciones comparadas y umbral final.

\texttt{\small .\allowbreak{}tolist(\allowbreak{})\allowbreak{}} convierte la tabla de NumPy en listas normales que JSON puede escribir. Conservar las opciones comparadas permite distinguir la elección por validación de la evaluación de prueba.


\begin{Codigo}{12B · Código}
informe.update({
    "modelo": "LogisticRegression · lbfgs · class_weight=balanced · píxeles/255",
    "C_elegido_solo_con_validacion": C_elegido,
    "opciones_C_y_resultados_train_val": resultados_C.tolist(),
    "umbral_predefinido": UMBRAL_FINAL,
})
print("Modelo registrado: C =", C_elegido, "· umbral =", UMBRAL_FINAL)
\end{Codigo}

\begin{Salida}{12B · Resultado}
Modelo registrado: C = 0.1 · umbral = 0.5
\end{Salida}

\needspace{10\baselineskip}\phantomsection\addcontentsline{toc}{subsection}{12C · Registrar evaluación y ejemplos}\Sub{12C · Registrar evaluación y ejemplos}\label{paso-12c}

Añadimos los resultados de prueba: matriz, medidas y primeros ejemplos de las casillas. \texttt{\small zip} empareja cada nombre de medida con su valor; \texttt{\small dict} convierte esas parejas en entradas.

La matriz se guarda como filas mediante \texttt{\small .\allowbreak{}tolist(\allowbreak{})\allowbreak{}}. Son las mismas cifras de los gráficos, sin elegir resultados diferentes para el archivo.


\begin{Codigo}{12C · Código}
informe.update({
    "matriz_prueba_filas_referencia_columnas_prediccion": matriz_test.tolist(),
    "metricas_prueba": dict(zip(["exactitud", "equilibrada", "sensibilidad", "especificidad"], metricas)),
    "metricas_referencia_prueba": base_test,
    "ejemplos_prueba": indices_ejemplos,
})
print("Evaluación registrada:", informe["matriz_prueba_filas_referencia_columnas_prediccion"])
\end{Codigo}

\begin{Salida}{12C · Resultado}
Evaluación registrada: [[158, 76], [14, 376]]
\end{Salida}

\needspace{10\baselineskip}\phantomsection\addcontentsline{toc}{subsection}{12D · Registrar versiones y límites}\Sub{12D · Registrar versiones y límites}\label{paso-12d}

Las versiones ayudan a entender por qué otra computadora podría obtener pequeñas diferencias. \texttt{\small \{p:\allowbreak{} version(\allowbreak{}p)\allowbreak{} for p in .\allowbreak{}.\allowbreak{}.\allowbreak{}\}} registra la versión de cada biblioteca y \texttt{\small sys.\allowbreak{}version.\allowbreak{}split(\allowbreak{})\allowbreak{}[\allowbreak{}0]\allowbreak{}} la de Python.

También guardamos lo que no comprobamos: independencia por paciente, calibración clínica y validación externa. Documentar un límite es parte de describir el experimento, no un resultado adicional del modelo.


\begin{Codigo}{12D · Código}
informe.update({
    "versiones": {p: version(p) for p in ["numpy", "matplotlib", "scikit-learn"]},
    "python": sys.version.split()[0],
    "limites": "Uso educativo; baja resolución; sin IDs para auditar independencia por paciente; sin validación clínica externa ni calibración clínica."
})
print("Versiones:", informe["versiones"])
\end{Codigo}

\begin{Salida}{12D · Resultado}
Versiones: {'numpy': '2.5.4', 'matplotlib': '3.11.2', 'scikit-learn': '1.9.1'}
\end{Salida}

\needspace{10\baselineskip}\phantomsection\addcontentsline{toc}{subsection}{12E · Guardar la ficha en JSON}\Sub{12E · Guardar la ficha en JSON}\label{paso-12e}

\texttt{\small json.\allowbreak{}dumps} transforma el diccionario en texto; \texttt{\small ensure\_\allowbreak{}ascii=\allowbreak{}False} mantiene acentos e \texttt{\small indent=\allowbreak{}2} lo ordena para leerlo. \texttt{\small write\_\allowbreak{}text(\allowbreak{}.\allowbreak{}.\allowbreak{}.\allowbreak{},\allowbreak{} encoding=\allowbreak{}"utf-8")\allowbreak{}} guarda \texttt{\small resultados\_\allowbreak{}modulo\_\allowbreak{}ml.\allowbreak{}json} en la carpeta actual.

En Colab lo encontrarás en el panel de archivos: actualiza la lista y descarga ese archivo si quieres conservar tu ejecución. El almacenamiento de la sesión es temporal; guardar el notebook no descarga automáticamente ese JSON.


\begin{Codigo}{12E · Código}
Path("resultados_modulo_ml.json").write_text(json.dumps(informe, ensure_ascii=False, indent=2), encoding="utf-8")
print("Guardado: resultados_modulo_ml.json")
\end{Codigo}

\begin{Salida}{12E · Resultado}
Guardado: resultados_modulo_ml.json
\end{Salida}

\needspace{10\baselineskip}\phantomsection\addcontentsline{toc}{subsection}{12F · Cerrar con el mapa del experimento}\Sub{12F · Cerrar con el mapa del experimento}\label{paso-12f}

Las tres tarjetas reúnen lo que hicimos: \textbf{aprender pesos}, \textbf{elegir con validación} y \textbf{evaluar con decisiones fijas}. No son tres entrenamientos independientes.

\texttt{\small axis(\allowbreak{}"off")\allowbreak{}} oculta ejes, las tuplas reúnen posición/texto/color y \texttt{\small for} recorre las tarjetas. \texttt{\small text} escribe cantidades calculadas y \texttt{\small bbox} dibuja el fondo. Para explicar tu resultado, indica los datos, C, umbral, matriz y límites; una cifra aislada no cuenta todo el experimento.


\begin{Codigo}{12F · Código}
fig, ax = plt.subplots(figsize=(11, 3.8), layout="constrained")
ax.axis("off")
tarjetas = [(0.16, "APRENDER", f"4690 imágenes\n784 características\nPesos ajustados", "#226b8a"), (0.50, "ELEGIR", f"516 imágenes\nC = {C_elegido:g}\nUmbral previo = 0.5", "#007f89"), (0.84, "EVALUAR", f"624 imágenes\nEquilibrada = {equilibrada:.1%}\nAciertos y errores visibles", "#c96932")]
for x, titulo, texto, color in tarjetas:
    ax.text(x, 0.55, titulo+"\n\n"+texto, ha="center", va="center", color="white", transform=ax.transAxes, bbox=dict(boxstyle="round,pad=1", facecolor=color, edgecolor="none"))
ax.text(0.5, 0.08, "Un resultado reproducible en este conjunto no es validación clínica.", ha="center", transform=ax.transAxes, color="#102c40")
plt.show()
print("Guardado: resultados_modulo_ml.json")
\end{Codigo}

\Figura{paso-12f-1.png}{12F · Aprender ajusta los pesos; elegir compara candidatos en validación; evaluar aplica las decisiones fijas a prueba.}{94mm}

\begin{Salida}{12F · Resultado}
Guardado: resultados_modulo_ml.json
\end{Salida}

\Sub{Comprueba que puedes explicar el experimento}

\begin{enumerate}
\item \textbf{¿Qué aprende \texttt{\small fit}?} Los pesos de las 784 entradas y el sesgo, usando imágenes y etiquetas de entrenamiento. No aprende el umbral al mover un control.
\item \textbf{¿Qué recibe \texttt{\small predict\_\allowbreak{}proba}?} Filas de intensidades. Aplica los pesos ya fijados y devuelve una puntuación por clase. No recibe la referencia de la imagen para adivinarla.
\item \textbf{Para la misma imagen, ¿qué cambia al bajar el umbral de 0.5 a 0.3?} La comparación y posiblemente la etiqueta predicha. Los píxeles, pesos y puntuación permanecen iguales.
\item \textbf{¿Por qué 73.84\% de exactitud puede ocultar un problema?} La regla mayoritaria acierta esa proporción en validación respondiendo siempre 1, pero falla todas las referencias 0; su exactitud equilibrada es 50\%.
\item \textbf{¿Podemos mirar prueba, cambiar C y presentar otra vez esa prueba como independiente?} No: ya habríamos usado su resultado para decidir. Debemos reconocer esa dependencia y evaluar nuevas decisiones con datos apropiados reservados.
\end{enumerate}


\Pagina{DEL EJERCICIO A UN ESTUDIO REAL}{Qué faltaría para ir más lejos}{Un buen resultado de taller es el comienzo de una pregunta.}
\begin{center}
\begin{tikzpicture}[node distance=7mm,box/.style={draw=Teal!40,fill=Pale,rounded corners=3pt,text width=145mm,inner sep=9pt,align=left,font=\sffamily},arr/.style={-{Stealth[length=2mm]},thick,Teal}]
\node[box] (a) {\textbf{Definir la pregunta}\quad ¿Qué entrada, población, decisión y uso se quiere estudiar?};
\node[box,below=of a] (b) {\textbf{Controlar la procedencia}\quad Permisos, etiquetas, calidad, identidad de pacientes y contexto de adquisición.};
\node[box,below=of b] (c) {\textbf{Separar correctamente}\quad Evitar que imágenes de la misma persona o copias entren en evaluación independiente.};
\node[box,below=of c] (d) {\textbf{Evaluar fuera del desarrollo}\quad Otros pacientes y, cuando corresponda, otros centros, equipos y períodos.};
\draw[arr] (a)--(b);\draw[arr] (b)--(c);\draw[arr] (c)--(d);
\end{tikzpicture}
\end{center}
\Sub{El cambio de contexto}
Un modelo puede funcionar en su conjunto de origen y fallar cuando cambia la población o la forma de adquirir las imágenes. Ese \textbf{cambio de distribución} puede afectar apariencias, etiquetas y frecuencia de clases. Un buen porcentaje en radiografías pediátricas reducidas no se transporta automáticamente a adultos ni a estudios clínicos de alta resolución.
\Sub{Las etiquetas también merecen revisión}
El aprendizaje supervisado aprende a aproximar las etiquetas que recibe. Si hay errores, sesgos o una definición ambigua de la tarea, puede reproducirlos. La evaluación respecto a una etiqueta publicada no demuestra por sí misma que esa etiqueta sea un estándar clínico perfecto.
\Sub{Antes de usar imágenes propias}
Usa una copia destinada a docencia y verifica los permisos correspondientes. Revisa cabeceras, información incrustada en píxeles y archivos acompañantes. Una conversión de formato no equivale a anonimización. El ejercicio utiliza exclusivamente las imágenes públicas indicadas.
\begin{Idea}{Un siguiente paso razonable}
Comparar una red convolucional con esta referencia, manteniendo el mismo diseño de separación y evaluación. Una arquitectura más compleja introduce otras decisiones; no corrige automáticamente los problemas de procedencia, etiquetas o dependencia entre imágenes.
\end{Idea}



\Pagina{FUENTES Y REPRODUCCIÓN}{De dónde viene este material}{Referencias consultadas y archivos para repetir el recorrido.}
\begin{enumerate}[label={[\arabic*]},leftmargin=2em]
\item \textbf{MedMNIST.} Descripción, alcance educativo y datos: \url{https://medmnist.com/}. Metadatos y huella del archivo: \href{https://github.com/MedMNIST/MedMNIST/blob/main/medmnist/info.py}{MedMNIST/info.py}. Descarga: \href{https://zenodo.org/records/10519652}{Zenodo, registro 10519652}.
\item \textbf{Yang, J. y colaboradores (2023).} \emph{MedMNIST v2: A large-scale lightweight benchmark for 2D and 3D biomedical image classification.} Scientific Data 10, 41.\\ \url{https://doi.org/10.1038/s41597-022-01721-8}.
\item \textbf{Kermany, D.; Zhang, K.; Goldbaum, M. (2018).} \emph{Labeled Optical Coherence Tomography (OCT) and Chest X-Ray Images for Classification.} Mendeley Data, versión 2. \url{https://doi.org/10.17632/rscbjbr9sj.2}.
\item \textbf{scikit-learn.} \href{https://scikit-learn.org/stable/modules/generated/sklearn.linear_model.LogisticRegression.html}{Regresión logística}; \href{https://scikit-learn.org/stable/modules/generated/sklearn.dummy.DummyClassifier.html}{referencia DummyClassifier}; \href{https://scikit-learn.org/stable/modules/model_evaluation.html}{métricas de evaluación}; \href{https://scikit-learn.org/stable/common_pitfalls.html}{separación de datos y errores frecuentes}.
\end{enumerate}
\Sub{Atribución y modificaciones}
PneumoniaMNIST y su fuente indicada se distribuyen bajo \href{https://creativecommons.org/licenses/by/4.0/}{Creative Commons Attribution 4.0 (CC BY 4.0)}. Se atribuye la colección derivada a Yang y colaboradores y las imágenes de origen a Kermany, Zhang y Goldbaum. MedMNIST ya realizó el recorte y la reducción a 28 × 28. Este taller conserva el archivo publicado, adapta la selección docente eliminando copias exactas mediante comparación de sus píxeles, normaliza intensidades y dibuja nuevas figuras. No hay respaldo implícito de los autores a este taller.
La portada y las figuras de resultados se componen a partir de las imágenes y cálculos del ejercicio; no son imágenes médicas generadas con IA. Los diagramas conceptuales son ilustraciones educativas. Las salidas son una ejecución de referencia, no una certificación diagnóstica.
\Sub{Recompilar el PDF}
El archivo \texttt{Fuente\_LaTeX\_modulo\_ML.zip} contiene texto, estilo, figuras, fuentes tipográficas y licencia. Extrae su contenido y compila \texttt{Guia\_modulo\_ML.tex} dos veces con \textbf{XeLaTeX}; conserva las carpetas junto al archivo principal.
\begin{Idea}{Identidad del archivo de datos}
\texttt{pneumoniamnist.npz}: 4 170 669 bytes.\\
MD5 publicado: \texttt{28209eda62fecd6e6a2d98b1501bb15f}.\\
Descarga y comprobación: 23 de noviembre de 2026. La huella SHA-256, el detalle de duplicados y las condiciones se documentan con los datos del módulo.
\end{Idea}
\end{document}
